\chapter{图像与视频模型}
\label{chap:vision-image-video-models}

一张$224\times224$的街景照片左侧有红色自行车，右侧有红色汽车；一段16帧的拿杯视频记录了一只手伸向桌上的杯子并把它拿起。对照片，我们可能要判断场景类别，也可能要沿轮廓分出汽车；对视频，我们可能要判断整段动作是不是“拿杯”，也可能要知道每一帧中手的位置和远近；我们还可能要求模型生成另一张街景，或续出拿杯之后的几帧。这些请求读取同样的像素，需要的表示却不同：类别只要一个整体向量，轮廓和深度需要保留位置的特征网格，生成则需要一种既紧凑、又能还原成像素的表示。

本章研究这些表示由什么计算形成、怎样选用。读者已经在前几卷学过卷积、残差连接、注意力和Transformer，也学过变分自编码器、扩散模型和自监督学习的目标构造；这里不重讲这些机制，只在用到时说明所需的结论和条件。第\ref{chap:vision-overview}章已经规定了观测、成果和约束，\ref{chap:vision-image-formation}章到\ref{chap:vision-video-foundations}章说明了像素怎样被采样和量化、图像怎样被滤波和分解、视频怎样按时间采样并估计运动。本章的许多结构可以看作这些经典运算的可学习版本：卷积层是一组由数据决定系数的空间滤波器，逐级降采样的特征图与图像金字塔的组织方式相同，图像块嵌入与分块变换编码一样对固定大小的块做线性变换，视频模型取帧时同样受时间采样的限制。我们在用到时点明这些对应关系，不重讲基础章的内容。

本章先把一个视觉模型拆成“输入的编码”和“输出的解码”两步，建立共同的形状记号。随后的三组模型分别回答三个问题：图像表示依靠什么结构和监督形成，视频表示怎样在可承受的计算下加入时间，生成表示怎样既压缩像素又能被解码还原。每个模型都按论文给出它的具体结构：卷积网络列出逐层或逐阶段的配置和输出尺寸，Transformer列出图像块大小、词元数、层数、宽度和头数，并写出单个模块内部的计算顺序。每组末尾在同一任务和同一输入下比较所介绍的模型，最后归纳这些模型在信息保留、训练数据和计算资源上的使用约束。后续任务章直接使用这里的模型，只说明任务中的用法；\ref{chap:vision-multimodal-models}章在这些视觉表示之上连接文字和声音。

\section{视觉模型的计算流程}
\label{sec:vision-ivm-pipeline}

图像的高宽因照片而异，视频的帧数因录像而异，而一个网络的参数规模是固定的；任务输出又可能是一个类别、一张与原图对齐的标签图，或者一段新画面。视觉模型因此要解决两个连续的问题：怎样把大小不定的像素数组组织成网络可以计算的单元，以及怎样把网络内部的表示读成任务需要的形状。前一步的难点在于既要控制单元数量，又要保留位置和时间；后一步的难点在于内部网格通常比原图粗，需要把粗网格上的判断还原到原图坐标。

两步都有不同路线。编码一侧的核心差别在于计算的基本单元：卷积网络以像素网格为单元，逐层用局部窗口汇总并降低分辨率，自然得到多个尺度的特征图；Transformer式编码器先把图像切成图像块、把视频切成时空块，再让所有单元互相读取，单元数量从一开始就由块的大小决定。解码一侧的差别在于保留多少空间结构：只读一个全局向量，读出每个位置的特征，或者把多个尺度的特征合起来恢复到像素分辨率。本节先说明编码，再说明解码；后面各节的模型都是这两步的具体实现。

\subsection{视觉输入的编码}
\label{sec:vision-ivm-encoding}

沿用式\eqref{eq:vision-image-array}的记号，图像为$\bm X\in\mathbb R^{h\times w\times c}$。编码器把它变成一组特征，每个特征对应原图的一个位置范围。卷积网络的做法是在像素网格上滑动局部窗口：一层步长为$s$的卷积把网格的高宽各缩小为原来的$1/s$，通道数则由卷积核个数决定，卷积的输出尺寸和感受野计算可参考\BookChapterRef{chap:convolutional-networks}{卷五的卷积神经网络章}。逐层累计后，第$\ell$个阶段的输出是特征图$\bm F^{(\ell)}\in\mathbb R^{(h/s_\ell)\times(w/s_\ell)\times d_\ell}$，$s_\ell$是该阶段相对原图的累计步长。以$224\times224$的街景为例，常见卷积主干在累计步长4、8、16、32处分别给出$56\times56$、$28\times28$、$14\times14$、$7\times7$的网格。

卷积层与\ref{sec:vision-transforms-spatial}节的空间滤波是同一种运算。均值、高斯和Sobel滤波器用一个固定的小核在图像上滑动，输出每个位置邻域像素的加权和；卷积层同样滑动小核，区别在于核的系数由训练决定，每层同时使用几十到几百个核，并在线性滤波之后接非线性激活。一层有64个输出通道的$3\times3$卷积，相当于并行运行64个空间滤波器，每个滤波器把它在所有输入通道上的响应相加。AlexNet第一层学到的核大多具有方向和频率选择性，另一些对颜色斑块敏感\scite{krizhevsky2012}，与人工设计的梯度算子和带通滤波器相似；更深的层在这些响应之上继续滤波，组合出纹理和部件。人工滤波器的作用可以从频率响应直接读出，学到的核没有这种保证，只能通过可视化或下游任务间接检查。

这几个尺度的特征承担不同职责。步长4的网格中一个位置只对应原图$4\times4$像素，适合定位细小边界，但每个位置看到的范围有限；步长32的网格中一个位置对应$32\times32$像素，经过多层汇总后能读到较大范围的上下文，却已经分不清这个格子里有几个小物体。多尺度特征（multi-scale features）就是同时保留这些不同步长的特征图，让后续预测在细节和上下文之间取用。街景右侧的汽车宽约80像素，在步长32的网格上只占两三个格子；左侧远处若有一块宽20像素的路牌，它在这一网格上还不到一个格子。

这组逐级减半的特征图与\ref{sec:vision-transforms-pyramid}节的高斯金字塔组织方式相同：金字塔的每一层先用高斯核平滑、再隔点取样，卷积网络的每个阶段也在若干层滤波之后以步长2降采样。两者的差别在于降采样之前是否保证低通。按采样定理，隔点取样会把高于新奈奎斯特频率的成分折叠成低频的假信号，所以金字塔先平滑再抽取；这一结论在一维情形下的推导，以及平均池化、最大池化与抽取的关系，可参考\BookChapterRef{chap:signal-analysis}{卷一的信号章}，二维图像中的混叠与摩尔纹见\ref{sec:vision-formation-aliasing}节。步长卷积的核由训练决定，最大池化取邻域最大值，二者都不保证先滤掉高频；一个可观察的后果是，输入平移一两个像素时，降采样后的特征可能明显变化。\ref{sec:vision-ivm-magvit2}节的MAGVIT-v2在判别器的下采样中加入模糊池化以增强平移不变性，就是把金字塔的抗混叠做法放回网络。金字塔的另一种用法也在网络中重新出现：\ref{chap:vision-image-recognition}章的特征金字塔网络把粗层的语义逐级上采样并与细层相加\scite{vision-1612-03144}，其计算顺序与拉普拉斯金字塔从粗到细逐层补回细节的重建相似。

Transformer式编码器换了一种组织方式。它先把图像切成互不重叠的\term{图像块}（image patch），也就是边长为$p$像素的小方块，每块作为一个计算单元。设$h,w$均可被$p$整除，则块数为$n=hw/p^2$。把第$i$块按固定顺序展开为$\bm x_i\in\mathbb R^{p^2c}$，用线性映射$\bm A\in\mathbb R^{d\times p^2c}$和偏置$\bm b\in\mathbb R^{d}$得到
\begin{equation}
  \bm h_i^{(0)}=\bm A\bm x_i+\bm b+\bm e_i,
  \qquad i=1,\ldots,n,
  \label{eq:vision-ivm-patch-embedding}
\end{equation}
其中$\bm e_i\in\mathbb R^{d}$是第$i$块的位置表示。对$224\times224\times3$的照片取$p=16$，得到$14\times14=196$个块，每块展开为$16\times16\times3=768$维，再映射为$d$维向量。这一步只做一次线性映射，块内部的空间结构没有经过逐层卷积；块与块之间的关系全部交给后续的注意力层。

切块之前，照片通常已经被调整到$224\times224$。这一步本身是一次重采样，所用的双线性或双三次插值，以及缩小时若不先平滑就会出现的混叠，见\ref{sec:vision-formation-interpolation}节；模型只能看到重采样之后的像素。图像块嵌入也可以与\ref{sec:vision-transforms-dct}节的分块离散余弦变换对照：两者都把图像切成不重叠的固定大小块，并对每一块施加同一个线性变换。区别在于变换矩阵的来源和用途。离散余弦变换的基是固定的正交基，可以精确逆变换，压缩来自对高频系数的粗量化；$\bm A$由训练得到，不要求可逆，768维的块展开映射到$d$维之后保留哪些成分，由训练目标决定。块的大小在两处起相同的作用：块越大，块内的空间关系越早被一次线性运算混合，后续计算也越少。

位置表示决定模型能否区分“左侧的车”和“右侧的车”。注意力本身对输入顺序不敏感，若去掉$\bm e_i$，把196个块的顺序打乱不会改变每个块读到的内容集合。位置表示可以是可学习的表格，也可以是按坐标计算的正弦函数或旋转位置编码，具体形式可参考\BookChapterRef{chap:transformer}{卷五的Transformer章}。可学习表格只为训练时的网格大小准备了向量；输入分辨率改变时，网格从$14\times14$变为$28\times28$，常见做法是对位置表在二维网格上插值。插值给出了新位置的向量，却不保证模型对更细的块划分仍然适用，这需要在新分辨率上验证或微调。按坐标计算的位置编码可以直接用于新网格，但训练中没有出现过的坐标范围同样需要检验。

视频还要处理时间。第一步是帧采样：从$\{(\bm X_t,\tau_t)\}$中按步长$k$选出$T$帧，模型实际看到的时间跨度约为$Tk/f$秒。拿杯视频若以每秒30帧录制、按步长4取16帧，输入覆盖约2.1秒；若同一段动作像例\ref{ex:vision-overview-sampling}那样落在两次采样之间，它就不在输入里。按\ref{sec:vision-videofound-frame-rate}节的时间采样分析，步长$k$把有效帧率降为$f/k$，频率高于$f/(2k)$的周期运动会发生时间混叠，例如转动的车轮看起来变慢甚至倒转；每一帧内部的运动模糊则由录制时的快门决定（\ref{sec:vision-videofound-shutter}节），取帧方式无法改变。视频文件通常按\ref{sec:vision-videofound-formats}节的YUV 4:2:0格式存储，送入模型前先解码并转换成RGB帧，色度分量在这一步已从四分之一的分辨率插值回来。第二步是切块。一种办法是逐帧切成图像块，再给每个块同时附加空间位置和帧序号；另一种办法是把连续$t$帧中同一位置的$p\times p$区域合成一个\term{时空块}（spatiotemporal patch），即覆盖$t\times p\times p$像素的立方体，再用线性映射嵌入。后者的单元数为
\begin{equation}
  n=\frac{T}{t}\cdot\frac{h}{p}\cdot\frac{w}{p}.
  \label{eq:vision-ivm-tube-count}
\end{equation}
16帧$224\times224$的拿杯视频取$t=2,p=16$，得到$8\times14\times14=1568$个单元；若逐帧切块，则有$16\times196=3136$个。时空块把单元数减半，代价是一个单元内部的两帧只经过一次线性混合，它的时间范围也不再是单个时间戳，而是两帧的区间。这次线性混合相当于对相邻两帧做一次可学习的时间滤波，并不估计运动：若手在两帧之间移动了半个块，块内混合的是两个位置的像素。运动关系要由后续注意力在多个块之间读出，这与\ref{sec:vision-videofound-block-matching}节的块匹配先搜索对应位置、再记录位移的做法不同。

编码之前的取舍已经限定了模型能看到的证据。设街景原图宽1200像素、高800像素，裁出左上角为$(400,200)$、宽高为$400\times400$的区域，再调整到$200\times200$作为局部输入。以连续边界坐标计，局部结果$(x',y')=(50,80)$对应原图$(x,y)=(400+2x',\,200+2y')=(500,360)$。输入记录因此至少要保存原图编号、裁剪边界和缩放关系；视频输入还要保存所取帧的真实时间戳，而不是让下游把采样后的帧当作相邻帧。从原始高分辨率照片重新裁取车窗，可以增加相对于缩小整图的可见细节；把已经缩小的车窗插值放大，只改变数组尺寸，不增加原始像素证据。

实际系统常把编码后的单元称为“视觉词元”。这里的词元只表示进入序列计算的一个单元，不一定来自离散词表；它与文字词元的整数编号、生成模型使用的离散码、原始像素值都不能直接互换。比较不同模型的词元数时，还要说明每个单元覆盖多大的空间和时间范围：同为196个单元，$p=16$的照片单元和$t=2,p=16$的视频单元承载的观测并不相同。

\subsection{任务输出的解码}
\label{sec:vision-ivm-decoding}

编码器交出的是一组带位置的特征，任务要的却是类别、标签图、深度图或者框。解码就是在这组特征上接一个任务头（task head），把特征读成任务输出的形状。解码方式取决于输出需要保留多少空间结构，常见的有三种。

第一种是全局向量。卷积网络通常对最后一个特征图做全局平均池化，把$7\times7\times d$的网格平均成一个$d$维向量；Transformer式编码器常在序列开头加一个专用的\term{分类位置}（class token, [CLS]），让它通过注意力汇总所有块，最后读取这一位置的输出，也可以对所有块的输出取平均。全局向量$\bm z\in\mathbb R^{d}$接一个线性层和softmax就得到类别概率。这种接口最便于分类和检索，但平均之后已经无法说出汽车在左边还是右边。

第二种是空间特征图。Transformer式编码器输出的196个块向量可以按原来的网格位置排回$14\times14\times d$的数组，卷积网络的中间特征图本身就是这种形式。在每个位置上接同一个线性层，就得到逐位置的预测，例如每个格子属于汽车、道路还是天空。这种接口保留了位置，但分辨率由编码步长决定：$14\times14$的网格中一个格子对应原图$16\times16$像素，汽车轮廓经过格子内部时，格子只能给出一个整体判断。

第三种是稠密预测头（dense prediction head），用于需要与原图同分辨率输出的任务，如语义分割和深度估计。\term{稠密预测}（dense prediction）指为每个像素或接近像素分辨率的每个位置都给出结果。稠密预测头把粗网格上的特征逐步上采样，并与较细尺度的特征相加或拼接，让上下文来自粗网格、边界位置来自细网格，最后输出$h\times w\times k$的数组。卷积主干天然提供多个尺度；Transformer式编码器只有一个$14\times14$的网格时，稠密预测头需要先用上采样和下采样从这一网格造出多个尺度，或者从编码器的几个中间层分别取特征。上采样本身可以用\ref{sec:vision-formation-interpolation}节的双线性插值，也可以用系数由训练决定的转置卷积。无论哪种做法，上采样都只能在已有特征之间插值和组合，编码时已经丢掉的细节不会因上采样而回来。

输出形状确定后，还要按第\ref{chap:vision-overview}章的要求把结果还原到原图坐标。逐位置预测的第$(i,j)$个格子对应原输入的$[jp,(j+1)p)\times[ip,(i+1)p)$像素范围，若输入经过裁剪和缩放，还要再按前面记录的变换回到原图。表\ref{tab:vision-ivm-output-forms}在同一张$224\times224$的街景上比较三种解码接口。

\begin{table}[htbp]
  \centering
  \small
  \begin{tabularx}{\linewidth}{@{}p{20mm}p{26mm}XX@{}}
    \toprule
    解码接口 & 输出形状 & 适合的任务 & 读不出的信息\\
    \midrule
    全局向量 & $d$ & 场景分类、整图检索 & 对象位于何处、有几个\\
    空间特征图 & $14\times14\times d$ & 粗粒度区域判断、区域检索 & 格子内部的边界\\
    稠密预测头 & $224\times224\times k$ & 分割、深度、关键点热图 & 编码时已丢失的细节\\
    \bottomrule
  \end{tabularx}
  \caption{同一街景输入下三种解码接口的比较，网格尺寸按$p=16$的块编码计算，为教学设定。后一行并不总是更好：输出越细，需要的标注和计算也越多。}
  \label{tab:vision-ivm-output-forms}
\end{table}
\FloatBarrier

\section{图像理解模型}
\label{sec:vision-ivm-image-understanding}

图像理解模型要把一张照片编码成可供多种任务读取的表示：分类需要整体向量能区分类别，分割和深度需要每个位置的特征能区分局部结构，检索需要相似的图像得到相近的向量。困难在于同一个表示要服务这些不同输出，而训练时通常只能使用一种监督。表示里保留什么，由两件事共同决定：网络结构预先假设了什么，例如邻近像素关系更紧密；训练目标要求它区分什么，例如一千个类别。

本节的三个模型对这两件事给出了不同回答。ResNet保留卷积的局部结构假设，用残差连接把网络做深，再用大规模类别标签训练，得到多尺度的特征图；ViT去掉了卷积的局部假设，把图像块当作序列交给Transformer，换来从第一层起就能跨全图交互的能力，但需要更多预训练数据来弥补缺少的结构假设；DINOv3沿用ViT的结构，改变的是监督来源，用不需要标签的自监督目标在十亿级图像上训练，并专门约束图像块特征的空间结构。三者依次改变了结构假设和监督来源，最后一个小节在同一街景任务上比较它们。

\subsection{ResNet}
\label{sec:vision-ivm-resnet}

卷积网络加深后理应能表示更复杂的函数，但直接堆叠层数会遇到训练困难。ResNet的原论文观察到，在CIFAR-10上，56层的普通卷积网络比20层的网络训练误差更高；这不是过拟合，因为问题出现在训练集上，而是优化器难以让更深的网络至少学到浅网络已有的映射\scite{he2016resnet}。卷五已经讲过ResNet的解决办法：残差连接把模块写成$\bm y=\bm x+\bm F_{\bm\theta}(\bm x)$，若某几层暂时不需要改变输入，残差分支只需接近零；反向传播时输入梯度多出一条经过恒等路径的项。本节不重复这些推导，只说明ResNet怎样作为视觉主干被使用。

\term{视觉主干}（visual backbone）指一个预先训练好的编码器，后续任务在它的输出上接任务头，而不重新设计编码部分。ResNet成为常用主干，一是因为它在ImageNet分类上训练稳定、可以做到上百层，二是因为它的四个阶段天然给出多尺度特征图。下面以ResNet-50为例，逐层说明一张$224\times224\times3$的照片怎样变成2048维向量。

网络的入口是一个“茎”（stem），负责尽快降低分辨率：conv1用64个$7\times7$、步长2的卷积核把图像变成$112\times112\times64$，随后一个$3\times3$、步长2的最大池化把网格缩到$56\times56$。从这里开始，四个阶段conv2\_x到conv5\_x分别堆叠3、4、6、3个瓶颈模块。每个阶段内部的网格大小不变，conv3\_1、conv4\_1和conv5\_1这三个阶段首块以步长2把高宽减半，同时把输出通道翻倍；论文只说明下采样由这三个首块完成，没有指明步长放在块内哪一个卷积上。最后，$7\times7\times2048$的特征图经全局平均池化成为2048维向量，接一千类的全连接层和softmax。表\ref{tab:vision-ivm-resnet50-layers}按论文的表1列出各层组的配置和输出尺寸。

\begin{table*}[tbp]
  \centering
  \small
  \begin{tabularx}{\linewidth}{@{}llXrrrr@{}}
    \toprule
    阶段或层组 & 层类型 & 卷积核与步长 & 输出通道 & 重复次数 & 输出尺寸 & 参数量\\
    \midrule
    输入 & — & — & 3 & — & $224\times224$ & —\\
    conv1 & 卷积 & $7\times7$，步长2 & 64 & 1 & $112\times112$ & 0.95万\\
    池化 & 最大池化 & $3\times3$，步长2 & 64 & 1 & $56\times56$ & 0\\
    conv2\_x & 瓶颈模块 & $1\times1{,}64\to3\times3{,}64\to1\times1{,}256$，步长1 & 256 & 3 & $56\times56$ & 21.6万\\
    conv3\_x & 瓶颈模块 & $1\times1{,}128\to3\times3{,}128\to1\times1{,}512$，首块步长2 & 512 & 4 & $28\times28$ & 122万\\
    conv4\_x & 瓶颈模块 & $1\times1{,}256\to3\times3{,}256\to1\times1{,}1024$，首块步长2 & 1024 & 6 & $14\times14$ & 710万\\
    conv5\_x & 瓶颈模块 & $1\times1{,}512\to3\times3{,}512\to1\times1{,}2048$，首块步长2 & 2048 & 3 & $7\times7$ & 1496万\\
    池化 & 全局平均池化 & $7\times7$ & 2048 & 1 & $1\times1$ & 0\\
    分类 & 全连接与softmax & — & 1000 & 1 & $1\times1$ & 205万\\
    \bottomrule
  \end{tabularx}
  \caption{ResNet-50的逐层结构，输入为$224\times224$的照片。卷积核、通道、重复次数和输出尺寸取自原论文表1；论文没有列出参数量，表中参数量按这些配置计算，包含批归一化的缩放与平移参数，阶段首块在通道变化时使用$1\times1$投影捷径（论文的选项B）。全网合计约2556万个参数，论文给出的计算量为$3.8\times10^9$次乘加。}
  \label{tab:vision-ivm-resnet50-layers}
\end{table*}

单个瓶颈模块内部的计算顺序是：$1\times1$卷积降维，$3\times3$卷积处理空间，$1\times1$卷积升维，再与捷径相加，最后经过ReLU。论文在每个卷积之后、激活之前使用批归一化（BN），第二个非线性放在相加之后。以conv2\_x中一个输入为$\bm x\in\mathbb R^{56\times56\times256}$的模块为例，
\begin{equation}
  \begin{aligned}
    \bm u_1&=\operatorname{ReLU}\bigl(\operatorname{BN}(\operatorname{Conv}_{1\times1}^{256\to64}(\bm x))\bigr)\in\mathbb R^{56\times56\times64},\\
    \bm u_2&=\operatorname{ReLU}\bigl(\operatorname{BN}(\operatorname{Conv}_{3\times3}^{64\to64}(\bm u_1))\bigr)\in\mathbb R^{56\times56\times64},\\
    \bm u_3&=\operatorname{BN}\bigl(\operatorname{Conv}_{1\times1}^{64\to256}(\bm u_2)\bigr)\in\mathbb R^{56\times56\times256},\\
    \bm y&=\operatorname{ReLU}\bigl(\bm u_3+\bm S(\bm x)\bigr),
  \end{aligned}
  \label{eq:vision-ivm-bottleneck}
\end{equation}
其中$\operatorname{Conv}_{k\times k}^{a\to b}$表示把$a$个通道变成$b$个通道的$k\times k$卷积。$\bm S$是捷径：输入与输出形状相同时取恒等映射，阶段首块的通道或分辨率改变时取一个$1\times1$卷积投影，必要时带步长2。图\ref{fig:vision-ivm-resnet-bottleneck}画出这一数据流。

\begin{figure}[htbp]
  \centering
  % ResNet-50 中 conv2_x 的一个瓶颈模块，结构依据 He 等（2016）表1与图5；宽约75mm。
\begin{tikzpicture}[font=\figurefont\small,
  io/.style={draw=FigureBlue,fill=FigureBlueFill,rounded corners=1.8pt,
    minimum width=34mm,minimum height=7mm,align=center,inner sep=1.2mm},
  conv/.style={draw=FigureYellow,fill=FigureYellowFill,rounded corners=1.8pt,
    minimum width=34mm,minimum height=7mm,align=center,inner sep=1.2mm},
  result/.style={io,draw=FigureRed,fill=FigureRedFill},
  flow/.style={-{Stealth[length=1.8mm]},draw=BookMuted,line width=0.8pt},
  note/.style={text=BookMuted,font=\figurefont\footnotesize,anchor=west}]
  \node[io] (x) at (0,0) {输入 $\bm x$};
  \node[conv] (c1) at (0,-1.25) {$1\times1$ 卷积，BN，ReLU};
  \node[conv] (c2) at (0,-2.5) {$3\times3$ 卷积，BN，ReLU};
  \node[conv] (c3) at (0,-3.75) {$1\times1$ 卷积，BN};
  \node[draw=BookMuted,circle,inner sep=0.6pt,line width=0.8pt] (add) at (0,-4.9) {$+$};
  \node[result] (y) at (0,-6.0) {ReLU，输出 $\bm y$};
  \draw[flow] (x) -- (c1);
  \draw[flow] (c1) -- (c2);
  \draw[flow] (c2) -- (c3);
  \draw[flow] (c3) -- (add);
  \draw[flow] (add) -- (y);
  \draw[flow] (x.east) -- ++(1.2,0) |- (add.east)
    node[pos=0.25,right,text=BookMuted,font=\figurefont\footnotesize] {捷径 $\bm S$};
  \node[note] at (-4.3,0) {$56\times56\times256$};
  \node[note] at (-4.3,-1.25) {$56\times56\times64$};
  \node[note] at (-4.3,-2.5) {$56\times56\times64$};
  \node[note] at (-4.3,-3.75) {$56\times56\times256$};
  \node[note] at (-4.3,-6.0) {$56\times56\times256$};
\end{tikzpicture}

  \caption{ResNet-50中conv2\_x的一个瓶颈模块。主路径依次为降维、空间卷积和升维，每个卷积之后接批归一化；右侧捷径在形状相同时为恒等映射，相加之后才经过最后一个ReLU。方框旁标出各步输出的高、宽和通道数。}
  \label{fig:vision-ivm-resnet-bottleneck}
\end{figure}

为什么要先降维再升维？把三层的参数数相加，$256\times64+9\times64\times64+64\times256=69632$；若直接在256个通道上做两层$3\times3$卷积，需要$2\times9\times256^2\approx118$万个参数，相差约17倍。论文的比较对象是ResNet-34的基本模块：在$56\times56$的网格上，两层64通道$3\times3$卷积有$2\times9\times64^2=73728$个参数，与瓶颈模块相近，而瓶颈模块输出的是256个通道。论文据此说明两种模块的计算量相近，并指出瓶颈设计主要出于训练时间的考虑。代价也来自这里：$3\times3$卷积只在64个通道上处理空间关系，信息必须先挤过这个较窄的中间层。捷径保持恒等映射同样是为了控制计算：瓶颈模块的捷径连接两端的高维特征，论文指出若把它换成投影，模块的计算量和参数量都会翻倍。

改变每个阶段的模块类型和模块数，就得到不同深度的ResNet。表\ref{tab:vision-ivm-resnet-family}列出论文表1的五种配置。18层和34层使用由两个$3\times3$卷积组成的基本模块，50层以上使用瓶颈模块；层数的计法是卷积层加最后的全连接层，例如ResNet-50为$1+3\times(3+4+6+3)+1=50$。

\begin{table*}[tbp]
  \centering
  \small
  \begin{tabular}{@{}lllrr@{}}
    \toprule
    模型 & 模块类型 & conv2\_x至conv5\_x的模块数 & 计算量 & 参数量\\
    \midrule
    ResNet-18 & 基本模块 & 2, 2, 2, 2 & $1.8\times10^9$ & 1169万\\
    ResNet-34 & 基本模块 & 3, 4, 6, 3 & $3.6\times10^9$ & 2180万\\
    ResNet-50 & 瓶颈模块 & 3, 4, 6, 3 & $3.8\times10^9$ & 2556万\\
    ResNet-101 & 瓶颈模块 & 3, 4, 23, 3 & $7.6\times10^9$ & 4455万\\
    ResNet-152 & 瓶颈模块 & 3, 8, 36, 3 & $11.3\times10^9$ & 6019万\\
    \bottomrule
  \end{tabular}
  \caption{ResNet系列在$224\times224$输入下的配置。模块数和计算量（乘加次数）取自原论文表1；参数量为按同一配置和投影捷径计算的结果，论文未直接给出。加深的模块主要位于conv4\_x，即$14\times14$这一尺度。}
  \label{tab:vision-ivm-resnet-family}
\end{table*}

从ResNet-50到ResNet-152，新增的34个模块中有30个放在conv4\_x，也就是$14\times14$的网格上；论文没有单独解释这一分配。ResNet-152的计算量为$11.3\times10^9$次乘加，仍低于VGG-16/19的$15.3/19.6\times10^9$次。

ResNet的训练使用ImageNet分类数据，约128万张训练图像、一千个类别。设第$i$张图像$\bm X_i$的类别为$y_i$，网络给出的类别概率为$p_{\bm\theta}(y\mid\bm X_i)$，训练最小化平均交叉熵
\begin{equation}
  L_{\mathrm{cls}}(\bm\theta)=-\frac{1}{n}\sum_{i=1}^{n}\log p_{\bm\theta}(y_i\mid\bm X_i).
  \label{eq:vision-ivm-cls-loss}
\end{equation}
这个目标只要求整图向量能区分一千个类别，并没有直接约束中间特征图上每个位置的含义。中间特征之所以能用于其他任务，是因为分类本身需要先识别边缘、纹理和物体部件；但哪些局部信息被保留、哪些被丢弃，仍由分类需要决定。论文的训练设置为：把图像短边随机缩放到256至480像素之间，随机裁出$224\times224$的区域并随机水平翻转；用批大小256的随机梯度下降，学习率从0.1开始，误差停止下降时除以10，最多训练$60\times10^4$次迭代，权重衰减$10^{-4}$，动量0.9，不使用dropout。测试时在比较实验中取10个裁剪平均，最好的结果则在多个尺度上以全卷积方式计算后平均。

这种逐层滤波再汇总的结构，可以与\ref{sec:vision-features-hog}节的HOG描述子对照。HOG先用梯度算子求每个像素的梯度，再在小格子里按方向统计直方图，最后在相邻格子组成的块内归一化；ResNet的前几层同样得到方向敏感的滤波响应，再经降采样把局部响应汇总到更粗的格子上。差别在于HOG的梯度算子、格子大小和方向分箱都由人事先选定，只有一层汇总；ResNet的滤波器和汇总方式都由分类目标决定，并叠加了几十层。HOG因此只描述边缘方向的局部分布，ResNet的深层还能表示部件和物体，但后者需要大量标注图像才能学出这些滤波器。

用作主干时有两种迁移方式。一种是冻结ResNet，只训练新接的任务头，所需标注少、计算省，但表示中缺少的信息无法补回；另一种是以ImageNet权重为初值，连同主干一起在新任务上微调，表示能向新任务调整，但需要更多标注，也可能在小数据上过拟合。检测和分割通常取后三或四个阶段的特征图，交给多尺度的任务头，相关用法在\ref{chap:vision-image-recognition}章展开。微调时还要注意批归一化：它在训练中按当前批次估计均值和方差，检测任务常因图像大、每批样本少而使这一估计不稳，工程上常把批归一化的统计量固定为预训练时的值。

沿街景例子看ResNet能提供什么。步长8的$28\times28$特征图仍能分开左侧自行车和右侧汽车所在的格子，步长32的$7\times7$特征图则让每个格子覆盖了较大范围的上下文，适合判断“这是一条街道”。但ResNet的每个位置只能通过逐层叠加的局部窗口间接读取远处内容，第一阶段的一个位置看不到画面另一侧的物体；它的特征也只受到类别监督，对没有类别标签区分的局部差别不一定敏感。

\subsection{ViT}
\label{sec:vision-ivm-vit}

ResNet的局部窗口在浅层只能读取邻近像素，远处内容要经过许多层才能汇合；而注意力允许任何两个位置在同一层直接交换信息。视觉Transformer（ViT）的想法是去掉卷积，把图像块组成的序列直接交给标准Transformer编码器\scite{dosovitskiy2020vit}。卷五已经介绍了它的输入构造和分类读出；这里关注它作为主干时输出什么、怎样训练和迁移，以及与ResNet相比得失在哪里。

ViT按式\eqref{eq:vision-ivm-patch-embedding}形成块向量，在序列开头加入一个可学习的分类位置，也就是一个不对应任何图像块的额外位置，专门用来汇总整幅图像。位置表示采用可学习的一维位置表：每个序号一个$d$维向量，与块向量相加。论文报告，换成显式利用二维坐标的位置表示没有带来明显提升；二维结构只在切块和改变分辨率时插值位置表这两处被人为引入。记第$i$块的展开为$\bm x_i\in\mathbb R^{p^2c}$，嵌入矩阵为$\bm E\in\mathbb R^{p^2c\times d}$，位置表为$\bm E_{\mathrm{pos}}\in\mathbb R^{(n+1)\times d}$，编码器的输入为
\begin{equation}
  \bm Z^{(0)}=\bigl[\bm x_{\mathrm{cls}};\ \bm x_1^\top\bm E;\ \ldots;\ \bm x_n^\top\bm E\bigr]+\bm E_{\mathrm{pos}}\in\mathbb R^{(n+1)\times d},
  \label{eq:vision-ivm-vit-input}
\end{equation}
它与式\eqref{eq:vision-ivm-patch-embedding}相同，只是改用行向量乘矩阵的写法并省略偏置，并在第一行放上分类位置$\bm x_{\mathrm{cls}}$。

编码器由$L$个相同的块组成。每块先做层归一化（LN）再进入子层，子层输出与输入相加，即“前置归一化”的残差结构：
\begin{equation}
  \begin{aligned}
    \bm Z'^{(\ell)}&=\bm Z^{(\ell-1)}+\operatorname{MSA}\bigl(\operatorname{LN}(\bm Z^{(\ell-1)})\bigr),\\
    \bm Z^{(\ell)}&=\bm Z'^{(\ell)}+\operatorname{MLP}\bigl(\operatorname{LN}(\bm Z'^{(\ell)})\bigr),
    \qquad \ell=1,\ldots,L,\\
    \bm y&=\operatorname{LN}\bigl(\bm Z^{(L)}_{0,:}\bigr).
  \end{aligned}
  \label{eq:vision-ivm-vit-block}
\end{equation}
$\operatorname{MSA}$是多头自注意力：$h$个头各自把$d$维向量投影到$d/h$维的查询、键和值，在全部$n+1$个位置之间做注意力，再把$h$个头的输出拼接并经一个$d\times d$的线性层投回$d$维。$\operatorname{MLP}$由两个全连接层和中间的GELU激活组成，先把$d$维扩到$4d$，再压回$d$。$\bm Z^{(L)}_{0,:}$是最后一层分类位置所在的行，归一化后得到整图表示$\bm y$。预训练时$\bm y$后接一个带一个隐藏层的MLP分类头，微调时换成一个零初始化的线性层。多头注意力和前馈层的一般定义，与位置编码一样在卷五的Transformer章中给出，这里只写出ViT的组织方式。

原论文按BERT的配置给出三个规模，表\ref{tab:vision-ivm-vit-configs}列出它们在$224\times224$输入下的形状。模型名中的“/16”表示块大小为16，例如ViT-L/16。

\begin{table*}[tbp]
  \centering
  \small
  \begin{tabular}{@{}lrrrrrrr@{}}
    \toprule
    模型 & 图像块大小 & 词元数 & 层数 & 宽度 & 头数 & MLP倍率 & 参数量\\
    \midrule
    ViT-B/16 & 16 & $196+1$ & 12 & 768 & 12 & 4 & 8600万\\
    ViT-L/16 & 16 & $196+1$ & 24 & 1024 & 16 & 4 & 3.07亿\\
    ViT-H/14 & 14 & $256+1$ & 32 & 1280 & 16 & 4 & 6.32亿\\
    \bottomrule
  \end{tabular}
  \caption{ViT的三个规模。层数、宽度、MLP宽度（3072、4096、5120，即宽度的4倍）、头数和参数量取自原论文表1；词元数按$224\times224$输入计算，“+1”为分类位置。块大小不改变参数量中的主体部分，只改变嵌入矩阵的输入维数和序列长度。}
  \label{tab:vision-ivm-vit-configs}
\end{table*}

以ViT-B/16为例走一遍形状。$224\times224$的照片切成196个$16\times16$块，每块展开为768维，乘以$\bm E$后仍为768维；加上分类位置共197行，每层输入输出都是$197\times768$。一个块的参数可以直接数出来：注意力的查询、键、值和输出投影共$4\times768^2$个权重，MLP有$2\times768\times3072$个权重，加上偏置和两个层归一化，约709万个参数。12个块约8505万，再加上嵌入矩阵、位置表和最后的层归一化，约8580万，与论文给出的8600万一致。每一层中，197个位置两两计算注意力权重，再经前馈层逐位置变换。输出端有两种可读的表示：分类位置的768维向量是整图表示，196个块的输出排回$14\times14\times768$就是空间特征图。与ResNet不同，ViT所有层都在同一个$14\times14$网格上计算，不会逐阶段降低分辨率，因此只有一个尺度；需要多尺度特征的稠密任务，要按\ref{sec:vision-ivm-decoding}节的做法另行构造。

ViT的训练沿用式\eqref{eq:vision-ivm-cls-loss}的类别监督，但对数据规模更敏感。原论文发现，只在ImageNet上训练时，较大的ViT不如同等规模的ResNet；在约1400万张图像、两万余类的ImageNet-21k，或3亿张图像的JFT-300M上预训练后再迁移，ViT才达到或超过ResNet。论文把这一现象解释为ViT缺少卷积所具有的局部性和平移等变性等归纳偏置，需要从更多数据中学出这些规律。这一解释与对比实验相符，但它只说明数据量和归纳偏置会共同影响结果，不能推出ViT在任何数据规模下都需要更多样本。

迁移时，ViT通常去掉预训练的分类头，换成新任务的线性层后整体微调。原论文还在比预训练更高的分辨率上微调：块大小保持16，$384\times384$的输入就有$24\times24=576$个块，预训练的位置表按二维网格插值到新尺寸。更多块带来更细的空间划分，也使注意力的计算按块数的平方增长：从196块到576块，每层需要计算的位置对从约3.8万增加到约33万。

在街景上，ViT的第一层就能让右侧汽车所在的块读取左侧自行车所在的块，适合需要全局关系的判断；但它的空间分辨率固定在16像素一格，宽20像素的路牌只占一两个块。ViT和ResNet一样只受类别标签监督，块特征是否保留了分类不需要的局部结构，同样没有直接保证。

\subsection{DINOv3}
\label{sec:vision-ivm-dinov3}

ResNet和ViT都依赖人工类别标签。标签限制了可用数据的规模，而且只监督整图类别，不直接要求块特征区分局部结构。自监督学习去掉了这一依赖：它从图像自身构造训练目标，例如让同一图像的不同裁剪得到一致的表示，或让被遮住的部分从其余部分被预测出来。卷八已经讲过这些目标的构造，包括DINO的学生—教师跨视图预测、教师的指数移动平均更新、居中与锐化，以及SwAV用Sinkhorn--Knopp迭代平衡原型分配的方法，可参考\BookChapterRef{chap:data-derived-knowledge-acquisition}{卷八“基于数据本身的知识获取”章}。DINOv3在这一路线上把模型和数据扩大到约67亿参数、十亿级图像，并面对一个扩大后才明显出现的问题：长时间训练后，整图表示仍在变好，块特征却在变差\scite{vision-2508-10104}。

论文报告，训练到约20万次迭代之后，ImageNet分类结果继续上升，而在VOC分割上用块特征做线性预测的结果开始下降，最大的模型最终降到早期水平以下。观察块特征可以看到原因的一部分：同一物体内部的块与块之间的相似度图变得嘈杂，分类位置与各块的余弦相似度随训练上升，说明块特征越来越像整图特征，失去了局部性。对分割和深度这类读取块特征的任务，这正是需要避免的变化。

DINOv3的主要想法是\term{Gram锚定}（Gram anchoring）：不直接要求块特征保持早期的数值，而是要求块与块之间的相似度结构保持早期教师的样子。把一张图的$P$个块特征按行排成矩阵并对每行做$L_2$归一化，记学生的为$\bm X_{\mathrm S}\in\mathbb R^{P\times d}$，作为参照的Gram教师的为$\bm X_{\mathrm G}$。$\bm X\bm X^\top$的第$(i,j)$元是第$i$块与第$j$块的余弦相似度，称为Gram矩阵。Gram锚定的损失为
\begin{equation}
  L_{\mathrm{Gram}}
  =\left\lVert\bm X_{\mathrm S}\bm X_{\mathrm S}^\top
  -\bm X_{\mathrm G}\bm X_{\mathrm G}^\top\right\rVert_{\mathrm F}^2.
  \label{eq:vision-ivm-gram}
\end{equation}
它只约束块之间的两两关系。若学生把所有块特征同时作同一个正交变换，Gram矩阵不变，损失也不变；因此特征本身仍可以随训练改进，只是“哪些块彼此相似”这一空间结构被锚定在块特征质量较好的早期状态。

用一个三块的教学例子算一次。设Gram教师的三个块特征为$(1,0)$、$(1,0)$、$(0,1)$，表示前两块属于同一物体；学生的为$(1,0)$、$(0.6,0.8)$、$(0,1)$，第二块已向第三块靠拢。两个Gram矩阵的非对角元分别为$(1,0,0)$和$(0.6,0,0.8)$，按$(1,2),(1,3),(2,3)$的次序列出。差值为$-0.4$、$0$、$0.8$，每个非对角差在矩阵中出现两次，对角元都为1，故损失为$2\times(0.16+0.64)=1.6$。梯度会把第二块拉回与第一块相似、与第三块不相似的结构。若学生把三个教师特征整体旋转$90^\circ$，得到$(0,1)$、$(0,1)$、$(-1,0)$，损失为零。

DINOv3的编码器是一个自行设计的ViT，记为ViT-7B，参数约67亿。它与式\eqref{eq:vision-ivm-vit-block}的块结构相同，仍是前置归一化的注意力子层加前馈子层，改动在规模和三个部件上。表\ref{tab:vision-ivm-dinov3-teacher}按论文表2列出它的配置，并与上一代DINOv2的教师对照。

\begin{table*}[tbp]
  \centering
  \small
  \begin{tabular}{@{}lll@{}}
    \toprule
    配置项 & DINOv2（ViT-giant） & DINOv3（ViT-7B）\\
    \midrule
    参数量 & 11亿 & 67亿\\
    层数 & 40 & 40\\
    图像块大小 & 14 & 16\\
    位置表示 & 可学习位置表 & 旋转位置编码\\
    寄存器数 & 4 & 4\\
    宽度 & 1536 & 4096\\
    前馈层 & SwiGLU，隐藏宽度4096 & SwiGLU，隐藏宽度8192\\
    头数（每头维数） & 24（64） & 32（128）\\
    DINO头（MLP各层宽度） & 4096--4096--256 & 8192--8192--512\\
    DINO原型数 & 12.8万 & 25.6万\\
    iBOT头（MLP各层宽度） & 4096--4096--256 & 8192--8192--384\\
    iBOT原型数 & 12.8万 & 9.6万\\
    \bottomrule
  \end{tabular}
  \caption{DINOv2与DINOv3教师网络的配置，数值取自DINOv3论文表2。两者层数相同，DINOv3加宽了每层并把块大小从14改为16。DINO头和iBOT头只在训练中使用，推理时丢弃。}
  \label{tab:vision-ivm-dinov3-teacher}
\end{table*}

三个部件各自对应一个问题。第一，前馈层使用SwiGLU：它用两路线性变换，一路经过SiLU激活后与另一路逐元素相乘，再经第三个线性层投回宽度，用门控代替ViT中的GELU前馈层；这种门控前馈层也在卷五的Transformer章中介绍过。第二，除分类位置外，序列中还加入4个\term{寄存器位置}（register token）：它们和分类位置一样不对应任何图像块，用来承接全局计算，减少块特征被挪作他用。第三，位置表示改为旋转位置编码，训练时还把坐标框$[-1,1]$随机缩放到$[-s,s]$，$s\in[0.5,2]$，使模型不依赖某一种固定的坐标尺度。以训练时$256\times256$的全局裁剪为例，输入为$16\times16=256$个块，加上1个分类位置和4个寄存器，序列长度为261，每层形状为$261\times4096$；$112\times112$的局部裁剪只有49个块。块大小从14改为16，使同样分辨率下的序列变短：DINOv2在$224\times224$下有256个块，DINOv3在$256\times256$下也是256个块。

训练分三个阶段。第一阶段是自监督预训练，目标为
\begin{equation}
  L_{\mathrm{pre}}=L_{\mathrm{DINO}}+L_{\mathrm{iBOT}}+0.1\,L_{\mathrm{KoLeo}}.
  \label{eq:vision-ivm-dinov3-pre}
\end{equation}
$L_{\mathrm{DINO}}$是整图层面的学生—教师目标：教师处理两个$256\times256$的全局裁剪，学生处理这两个全局裁剪和8个$112\times112$的局部裁剪，学生的分类位置输出要在一组原型上的分布与教师对另一视图的分布一致。$L_{\mathrm{iBOT}}$把同样的思路用到块层面：学生输入中一部分块被替换为掩码向量，学生要在这些位置预测教师看到完整图像时在同一位置给出的原型分布，这一项直接监督块特征。两项中教师分布原来用居中来防止所有样本集中到同一原型，DINOv3改用Sinkhorn--Knopp迭代平衡批内的原型分配。$L_{\mathrm{KoLeo}}$让一批中归一化后的整图特征彼此散开：设$\bm z_1,\ldots,\bm z_m$是批内特征，$\rho_i=\min_{j\neq i}\lVert\bm z_i-\bm z_j\rVert_2$是第$i$个特征到最近邻的距离，则$L_{\mathrm{KoLeo}}=-\frac{1}{m}\sum_{i}\log\rho_i$。若两个样本的特征几乎重合，$\rho_i$接近零，$-\log\rho_i$很大，梯度会把它们推开。论文在每16个样本的小批内计算这一项。

预训练使用论文整理的LVD-1689M数据集，约16.89亿张图像，从约170亿张Instagram公开图像中经聚类均衡和检索筛选而来，另混入ImageNet等数据；共训练100万次迭代，学习率、权重衰减和教师的移动平均系数都保持常数，以便在任何时刻都可以停止或进入下一阶段。第二阶段是Gram锚定精调：在100万次迭代之后加入式\eqref{eq:vision-ivm-gram}，只作用于全局裁剪，总目标为$L_{\mathrm{ref}}=w_{\mathrm D}L_{\mathrm{DINO}}+L_{\mathrm{iBOT}}+w_{\mathrm{DK}}L_{\mathrm{KoLeo}}+w_{\mathrm{Gram}}L_{\mathrm{Gram}}$。Gram教师取训练早期的教师网络，此后每1万次迭代更新为当前的教师。论文还让Gram教师处理两倍分辨率的图像，再把它的块特征图用双三次插值缩小一半，与学生的网格对齐后计算Gram矩阵；较高分辨率下得到的特征图更平滑，提供了更干净的相似度结构。第三阶段是高分辨率适应：再训练1万次迭代，全局裁剪取512或768像素，局部裁剪取112至336像素，并继续使用Gram锚定。论文报告，去掉这一项时高分辨率适应后的结果明显变差。

ViT-7B推理代价很高，论文又以冻结的ViT-7B为教师，把它蒸馏为一族较小的模型：标准的ViT-S、ViT-B、ViT-L，自行设计的ViT-S+和ViT-H+，以及四个ConvNeXt卷积模型。蒸馏沿用第一阶段的训练目标，但教师固定为ViT-7B本身，不再是学生的移动平均；论文没有在学生上观察到块特征退化，因此不使用Gram锚定。学生先训练100万次迭代，再以余弦日程降低学习率25万次迭代，最后做不带Gram锚定的高分辨率适应。表\ref{tab:vision-ivm-dinov3-family}列出这一族模型的规模和计算量。

\begin{table*}[tbp]
  \centering
  \small
  \begin{tabular}{@{}lrrr@{}}
    \toprule
    模型 & 参数量 & 推理计算量（$256^2$） & 推理计算量（$512^2$）\\
    \midrule
    ViT-S & 2100万 & 12 & 63\\
    ViT-S+ & 2900万 & 16 & 79\\
    ViT-B & 8600万 & 47 & 216\\
    ViT-L & 3亿 & 163 & 721\\
    ViT-H+ & 8.4亿 & 450 & 1903\\
    ViT-7B & 67.16亿 & 3550 & 14515\\
    ConvNeXt-T & 2900万 & 5 & 20\\
    ConvNeXt-S & 5000万 & 11 & 46\\
    ConvNeXt-B & 8900万 & 20 & 81\\
    ConvNeXt-L & 1.98亿 & 38 & 152\\
    \bottomrule
  \end{tabular}
  \caption{DINOv3模型族的参数量与推理计算量（单位GFLOPs），取自论文表14。ViT学生的块大小均为16；论文只给出ViT-S+和ViT-H+的参数量，没有给出它们的层数和宽度。ConvNeXt是卷积网络，没有分类位置，计算量基本随像素数线性增长。}
  \label{tab:vision-ivm-dinov3-family}
\end{table*}

分辨率从256加倍到512时，ViT模型的计算量增加到原来的4.1至5.3倍：块数变为4倍，前馈层和投影的计算随之变为4倍，注意力的位置对则变为16倍，模型越小，注意力在总计算中占比越大，倍数也越高；ConvNeXt的计算量基本随像素数线性增长，约为4倍。论文报告，最大的学生ViT-H+参数约为教师的八分之一，在其评测中与教师接近。此外，论文还按冻结视觉编码器、从头训练文字编码器的方式得到一个与文字对齐的版本，这一类图文对齐的机制在\ref{chap:vision-multimodal-models}章讨论。

使用DINOv3时，常见做法是冻结编码器，只训练很轻的任务头：分类用分类位置向量接线性层，分割和深度用$h/16\times w/16$的块特征图接线性层或稠密预测头。冻结的好处在于同一份特征可以供多个任务共用，DINOv3的块特征也正是为这种用法而训练的；但冻结同时意味着编码器没有见过的类别差别、任务特有的细节都无法由任务头补出。\ref{chap:vision-image-recognition}章“类别与场景的判定”给出分类头的训练与使用。

\subsection{图像表示的形成与接口选择}
\label{sec:vision-ivm-image-interface}

三个模型可以放在同一个问题上比较：用$224\times224$的街景照片，既要判断场景类别，又要分出右侧汽车的像素。表\ref{tab:vision-ivm-image-models}列出它们在结构假设、监督来源和输出接口上的差别。

\begin{table}[htbp]
  \centering
  \small
  \begin{tabularx}{\linewidth}{@{}p{16mm}XXX@{}}
    \toprule
    比较项 & ResNet & ViT & DINOv3\\
    \midrule
    结构假设 & 局部卷积，逐阶段降分辨率 & 图像块序列，全局注意力 & 同ViT，另加寄存器与旋转位置编码\\
    监督来源 & ImageNet类别标签 & 大规模类别标签预训练 & 无标签图像的自监督目标\\
    整图接口 & 池化后的向量 & 分类位置向量 & 分类位置向量\\
    空间接口 & 步长4至32的四个尺度 & 步长16的单一网格 & 步长16的单一网格，块特征受直接监督\\
    代表规模 & ResNet-50，2556万参数 & ViT-B/16，8600万参数 & ViT-7B约67亿参数，蒸馏学生2100万至8.4亿参数\\
    主要代价 & 远距离关系需多层汇合 & 依赖大规模预训练，计算随块数平方增长 & 训练计算与数据整理规模大，需蒸馏以供使用\\
    \bottomrule
  \end{tabularx}
  \caption{同一街景任务下三个图像理解模型的比较。表中只列结构与训练上的差别，各模型在具体基准上的成绩以相应任务章及原论文为准。}
  \label{tab:vision-ivm-image-models}
\end{table}
\FloatBarrier

ViT改变的是ResNet的结构假设：用图像块和全局注意力代替局部卷积，换来浅层即可读取全图关系，也使模型规模和数据规模更容易一同扩大；付出的代价是对预训练数据量更敏感，并失去了卷积主干天然的多尺度输出。DINOv3保留ViT的结构，改变的是监督来源：自监督目标不需要标签，可以使用规模大得多的图像，iBOT项和Gram锚定又直接约束块特征，使冻结的块特征适合分割和深度；付出的代价是训练所需的数据整理和计算规模大得多，并且需要蒸馏才能得到可用的小模型。

接口选择因此取决于任务读的是哪一种表示。只需判断场景类别时，三者的整图向量都可用，差别在于预训练覆盖的类别和视觉变化；需要分出汽车像素时，ResNet的多尺度特征图便于构造稠密预测头，DINOv3的块特征在冻结时就保留了较好的局部结构，而以分类为目标的ViT块特征通常需要连同主干一起微调。三者的联系也很直接：ViT的分类位置和块序列是DINOv3的计算骨架，ResNet的多尺度思想则在ViT之上以稠密预测头的形式重新出现；DINOv3蒸馏出的ConvNeXt模型又把自监督得到的表示传回卷积结构。

三者的空间接口也可以放回第一部分的经典方法中理解。ResNet各阶段的特征图对应一座可学习的图像金字塔，粗层负责上下文、细层负责定位，这与\ref{sec:vision-transforms-pyramid}节用不同层级分别处理大小结构的思路相同。DINOv3的Gram矩阵则是块与块之间的相似度矩阵，\ref{sec:vision-features-graph-segmentation}节的归一化割正是在这样一张相似度图上把像素分成几组。Gram锚定保住的就是这种分组所依赖的结构：同一辆汽车的块彼此相似、与道路的块不相似。差别在于，经典方法的相似度由颜色、纹理等人工特征定义，DINOv3的相似度由自监督训练得到，不限于颜色和纹理是否相近；代价是哪些部位被视为相似不再能由人直接规定，只能在具体图像上检查。

\section{视频理解模型}
\label{sec:vision-ivm-video-understanding}

视频理解模型要把一段帧序列编码成表示，供动作分类、时序定位和逐帧的局部预测使用。与图像相比，视频新增了两个困难。一是计算量：16帧视频即使用时空块也有1568个单元，若所有单元两两做注意力，每层要计算约246万个位置对，是单张图像的64倍。二是监督：逐段标注动作的成本远高于图像分类，而相邻帧高度相似，许多标注实际上在重复相同的画面。

本节的三个模型对这两个困难给出不同回答。TimeSformer保留有标签的监督训练，改变注意力的结构，把时空注意力分解为时间和空间两步，以降低计算量；VideoMAE改变监督来源，遮住九成时空块，让模型从剩下的部分重建像素，从无标签视频中学习，同时只编码可见的少数单元；V-JEPA 2.1沿用遮挡预测，但把预测目标从像素换成教师编码器的表示，并让可见位置和多个中间层也受到监督，使表示适合逐位置的密集预测。三者依次改变了注意力结构、监督来源和预测目标。

\subsection{TimeSformer}
\label{sec:vision-ivm-timesformer}

在TimeSformer之前，视频分类的主流做法是三维卷积网络，它用覆盖几帧、几个像素的局部卷积核同时处理时间和空间，远距离的时间关系要经过多层才能汇合。把ViT直接推广到视频，让所有时空块两两做注意力，可以在一层内读取任意帧的任意位置，但计算量按块数平方增长。TimeSformer要回答的问题是：能否只用注意力、不用卷积做视频分类，同时把计算控制在可接受的范围内\scite{vision-timesformer2021}。

它的主要想法是\term{分解时空注意力}（divided space-time attention）：每一层先让每个块只与其他帧中同一空间位置的块做注意力，再让它只与同一帧内的其他块做注意力。时间注意力负责“同一位置随时间怎样变化”，空间注意力负责“同一时刻画面各处的关系”；两步连用后，一个块经两次读取就能间接接触到任意帧的任意位置。

设输入为$F$帧、每帧切成$N$个块，$\bm z_{p,t}\in\mathbb R^{d}$是第$t$帧第$p$个块在某一层的表示，另有一个分类位置。一层的计算为
\begin{equation}
  \begin{aligned}
    \bm z'_{p,t}&=\bm z_{p,t}+\operatorname{Attn}_{\mathrm{time}}
      \bigl(\bm z_{p,t};\{\bm z_{p,t'}\}_{t'=1}^{F}\bigr),\\
    \bm z''_{p,t}&=\bm z'_{p,t}+\operatorname{Attn}_{\mathrm{space}}
      \bigl(\bm z'_{p,t};\{\bm z'_{p',t}\}_{p'=1}^{N}\bigr),\\
    \bm z^{\mathrm{out}}_{p,t}&=\bm z''_{p,t}+\operatorname{MLP}\bigl(\operatorname{LN}(\bm z''_{p,t})\bigr).
  \end{aligned}
  \label{eq:vision-ivm-divided-attention}
\end{equation}
$\operatorname{Attn}(\bm u;S)$表示以$\bm u$为查询、在集合$S$及分类位置上读取的多头注意力，读取前先做层归一化；时间注意力和空间注意力使用各自的查询、键、值投影矩阵。按每个查询需要比较的键计数，联合时空注意力要比较$NF+1$个，分解注意力只需$N+F+2$个。默认输入为8帧$224\times224$、块大小16，$N=196,F=8$，前者为1569个，后者为206个，相差约7.6倍。

把式\eqref{eq:vision-ivm-divided-attention}展开成一个块内部的计算顺序：先对输入做层归一化，以同一空间位置的8个块和分类位置为键值做时间注意力，经输出投影后加回输入；再对结果做层归一化，换一组投影矩阵，以同一帧的196个块和分类位置为键值做空间注意力，同样加回；最后经层归一化、MLP和残差相加得到本块输出。论文说明，时间注意力的输出不先经过MLP，而是直接作为空间注意力的输入。

模型输入是从视频中以1/32的采样率取出的8帧，即每32帧取一帧；每帧切成$16\times16$的块，第$t$帧第$p$块的展开$\bm x_{p,t}\in\mathbb R^{768}$经线性映射后加上可学习的时空位置表示$\bm e^{\mathrm{pos}}_{p,t}$，即每个“帧序号—空间位置”组合各有一个向量。编码器沿用ViT-B的12层、768宽度和12个头，每层按式\eqref{eq:vision-ivm-divided-attention}计算，形状保持$(8\times196+1)\times768$。最后对分类位置的输出做层归一化，接一个带一个隐藏层的MLP，给出动作类别概率。表\ref{tab:vision-ivm-timesformer-configs}列出论文中的三种输入配置。

\begin{table*}[tbp]
  \centering
  \small
  \begin{tabular}{@{}lrrrrrrr@{}}
    \toprule
    模型 & 输入 & 图像块大小 & 词元数 & 层数 & 宽度 & 头数 & 参数量\\
    \midrule
    TimeSformer & $8\times224^2$ & 16 & $8\times196+1$ & 12 & 768 & 12 & 1.214亿\\
    TimeSformer-HR & $16\times448^2$ & 16 & $16\times784+1$ & 12 & 768 & 12 & 论文未单列\\
    TimeSformer-L & $96\times224^2$ & 16 & $96\times196+1$ & 12 & 768 & 12 & 论文未单列\\
    \bottomrule
  \end{tabular}
  \caption{TimeSformer的三种配置，编码器均为ViT-B结构，MLP倍率为4。输入与块大小取自原论文，词元数按“帧数$\times$每帧块数$+1$”计算。分解注意力版本有1.214亿个参数，比同一结构的联合时空注意力版本（8590万）多出的部分来自每层独立的时间注意力参数；论文给出的单个视图推理计算量为0.59 TFLOPs。}
  \label{tab:vision-ivm-timesformer-configs}
\end{table*}

TimeSformer以动作类别标签监督训练，主要数据为Kinetics-400，约24万段训练视频、400个动作类别，目标与式\eqref{eq:vision-ivm-cls-loss}相同。参数初始化依赖图像预训练：每一块中的时间注意力和空间注意力都用ImageNet预训练ViT对应注意力层的同一组权重初始化。训练从完整视频中按1/32的帧率随机取片段，把短边随机缩放后裁出$224\times224$，批大小16，在32个GPU上用动量0.9、权重衰减$10^{-4}$的同步随机梯度下降。论文报告，在Kinetics-400上，用ImageNet-21k初始化得到78.0\%的准确率，用ImageNet-1k初始化为75.8\%，从零开始训练仅为64.8\%。推理时只从视频中部取一段片段，把短边缩放到224像素后取3个空间裁剪，平均3次的softmax得分。论文还给出两个变体：TimeSformer-HR使用16帧$448\times448$的输入以提高空间分辨率，TimeSformer-L使用96帧$224\times224$、每4帧取一帧的输入，单段覆盖约12秒。

在拿杯视频上，时间注意力让杯子所在位置的块能比较8个时刻的外观，看出杯子被抬起；空间注意力再把这一变化与手所在的块联系起来。但这种分解也有前提：时间注意力只比较同一空间位置，若杯子移动跨越了好几个块，同一位置在不同帧中已不是同一物体，模型需要依靠空间注意力在另一步中把它们联系起来。这一前提可以与\ref{sec:vision-videofound-motion}节的运动估计对照：块匹配在一个搜索窗口内寻找当前块在下一帧的对应位置，Lucas--Kanade假设邻域内运动很小并用梯度求位移；TimeSformer的时间注意力相当于把搜索窗口缩到同一个块，对运动的处理完全交给后续的空间读取，它既不输出位移，也不保证同一物体被对齐。位置表示的消融也说明时间信息依赖显式编码：去掉时间位置、只保留空间位置时，Kinetics-400上的准确率从78.0\%降为77.8\%，Something-Something V2上则从59.5\%降为52.5\%，后者的动作类别更依赖先后顺序。TimeSformer仍需要动作标签和图像预训练权重，标签数量直接限制它能学到的动作范围。

\subsection{VideoMAE}
\label{sec:vision-ivm-videomae}

TimeSformer依赖动作标签和图像预训练。卷八介绍的掩码自编码器（MAE）表明，图像可以通过“遮住一部分块、重建被遮的像素”进行自监督预训练，常用的遮挡比例为75\%。把这一做法直接搬到视频会遇到一个视频特有的问题：相邻帧几乎相同，若每帧独立随机遮挡，某帧被遮住的块在相邻帧同一位置往往可见，模型只要复制邻帧的像素就能完成重建，学不到关于运动和物体的知识。这种时间冗余正是视频编码所利用的：\ref{sec:vision-videofound-block-matching}节的块匹配在前一帧中找到最相似的块，只需记录位移和残差，就能用很少的比特表示当前帧。对视频编码有利的冗余，对掩码重建却是捷径。VideoMAE针对这一点设计遮挡方式\scite{vision-videomae2022}。

它的主要想法是\term{管状掩码}（tube masking）配合很高的遮挡比例：在空间网格上随机选出要遮的位置，这些位置在所有时间步都被遮住，像一根贯穿时间的管子。这样被遮的位置在任何一帧都不可见，无法从邻帧复制，模型只能依靠周围可见区域和整段的运动规律来推断。视频的时间冗余又使得遮掉90\%后仍然存在足够的信息可以推断，因此遮挡比例可以比图像高得多。

模型输入为16帧$224\times224$的片段，在Kinetics-400上每4帧取一帧，在Something-Something V2上每2帧取一帧。按式\eqref{eq:vision-ivm-tube-count}取$t=2,p=16$，切成$8\times14\times14=1568$个时空块，每块包含$2\times16\times16\times3=1536$个像素值，线性嵌入后加位置表示。管状掩码遮掉90\%，编码器只处理剩下约157个可见块；编码器是ViT-B，12层、768维，在可见块之间做联合时空注意力。解码器较轻，4层、384维：它把编码器的输出放回原位置，在被遮位置填入同一个可学习的掩码向量，经注意力计算后，为每个被遮块输出1536个像素值。表\ref{tab:vision-ivm-videomae-layers}按论文附录表8列出各层的形状，记法为“通道$\times$时间$\times$空间”。

\begin{table*}[tbp]
  \centering
  \small
  \begin{tabularx}{\linewidth}{@{}llXrl@{}}
    \toprule
    阶段 & 层类型 & 配置 & 重复次数 & 输出尺寸\\
    \midrule
    数据 & 帧采样 & Kinetics-400步长4，Something-Something V2步长2 & — & $3\times16\times224^2$\\
    时空块嵌入 & 三维卷积 & 核$2\times16\times16$，步长$2\times16\times16$，768通道 & 1 & $768\times8\times196$\\
    掩码 & 管状掩码 & 遮挡比例$\rho=0.9$，各帧遮同一组空间位置 & — & $768\times8\times[196(1-\rho)]$\\
    编码器 & Transformer块 & 联合时空注意力768宽，MLP宽3072 & 12 & $768\times8\times[196(1-\rho)]$\\
    投影 & 线性层 & $768\to384$，在被遮位置拼入可学习掩码向量 & 1 & $384\times8\times196$\\
    解码器 & Transformer块 & 联合时空注意力384宽，MLP宽1536 & 4 & $384\times8\times196$\\
    输出投影 & 线性层 & $384\to1536$ & 1 & $1536\times8\times196$\\
    重排 & — & 每个1536维向量排回$3\times2\times16\times16$ & — & $3\times16\times224^2$\\
    \bottomrule
  \end{tabularx}
  \caption{以16帧ViT-B为例的VideoMAE结构，取自原论文附录表8。时空块嵌入等价于对$2\times16\times16$的不重叠立方体做线性映射；掩码之后每帧保留约19.6个、全段约157个位置。解码器宽度为编码器的一半，预训练结束后丢弃。}
  \label{tab:vision-ivm-videomae-layers}
\end{table*}

解码器的深度在视频上比在图像上更重要。论文的消融显示，4个块的解码器结果最好；而在图像MAE中，较浅的解码器已经够用。论文没有给出这一差别的机制解释，我们只能确认它在其设置下成立。编码器之所以能只处理一成的位置，是因为管状掩码去掉的位置从输入序列中整个删除，注意力的位置对随之从约246万降到约2.5万。

训练目标只在被遮块上计算。记被遮块的索引集合为$\Omega$，第$i$块按块内均值和标准差归一化后的像素为$\bm x_i\in\mathbb R^{1536}$，解码器的输出为$\hat{\bm x}_i$，则
\begin{equation}
  L_{\mathrm{VideoMAE}}=\frac{1}{|\Omega|}\sum_{i\in\Omega}
  \lVert\hat{\bm x}_i-\bm x_i\rVert_2^2.
  \label{eq:vision-ivm-videomae}
\end{equation}
只在被遮位置计算损失，是因为可见块的像素已经输入编码器，重建它们不需要推断。预训练只用无标签视频，在Something-Something V2和Kinetics-400上默认训练800个轮次。预训练结束后丢掉解码器，把全部1568个块送入编码器，在其输出上接分类头，用动作标签微调整个编码器；推理时同样不再遮挡。微调后的ViT-B编码器有8700万个参数，论文给出的每个16帧视图的计算量约为180 GFLOPs，在Something-Something V2上测试时取2段$\times$3个裁剪共6个视图平均。

论文的消融说明了管状掩码的作用。在Something-Something V2上，以ViT-B、16帧、800轮预训练、相同微调协议比较，90\%的管状掩码得到69.6\%的准确率，同比例的逐块随机掩码为68.3\%，按整帧遮挡为61.5\%；遮挡比例取75\%时，管状掩码为68.0\%。同一设置下，从零训练为32.6\%，用ImageNet-21k有监督预训练初始化为61.8\%。论文还报告，预训练数据与目标数据之间的领域差别会影响迁移结果，数据质量比数量更重要。这些结论都依赖它所用的数据集和训练日程。

在拿杯视频上，管状掩码可能遮住杯子所在的整列位置，模型必须从手的运动方向和桌面周围的可见区域推断杯子在各帧的样子。VideoMAE的代价也来自这一目标：重建像素要求编码器保留纹理、光照等细节，其中有些细节本身难以从可见部分预测，模型的一部分容量被用来拟合这些内容。

\subsection{V-JEPA 2.1}
\label{sec:vision-ivm-vjepa}

\term{联合嵌入预测架构}（joint-embedding predictive architecture, JEPA）对VideoMAE的像素目标作了改变：从可见部分预测被遮部分在某个编码器中的表示，而不是被遮部分的像素。表示可以省去难以预测的像素细节，只保留可以从上下文推断的内容。早期的视频JEPA只在被遮位置计算预测误差，可见位置的输出没有直接目标，容易变成替其他位置汇总全局信息的通道，不利于逐位置的密集预测。V-JEPA 2.1针对这一点改进训练目标，要求同一个编码器既保持整体理解，又给出空间结构清楚、时间上连贯的密集特征\scite{vision-vjepa21-2026}。

密集表示是指每个时空网格位置都有一个向量，而不是每个位置已经有正确的语义或几何答案。V-JEPA 2.1的主要想法有三点：预测误差同时作用于被遮位置和可见位置；在编码器的多个中间层上都施加自监督；用二维和三维两种切块投影让图像和视频共享同一个主体编码器。前两点直接针对可见位置特征缺少约束的问题，第三点让大量图像数据也能参与训练。

设视频切块后的索引集合为$I$，可见位置集合为$C$、遮蔽位置集合为$M$，满足$I=C\cup M$且两者不相交。每个位置包含空间网格坐标和时间编号。学生编码器$E_{\bm\theta}$只处理可见块；预测器$P_{\bm\phi}$接收可见特征与携带缺失位置的遮蔽向量，为每个位置输出$\widehat{\bm{h}}_i\in\mathbb{R}^{d}$。教师编码器$E_{\bar{\bm\theta}}$处理同一次增强后的完整输入，给出目标$\bm{h}^{*}_i\in\mathbb{R}^{d}$。先固定两端所对应的位置，再比较向量差，避免把某处的手部特征拿去监督另一处的杯身。为了同时约束缺失与可见位置，一条训练样本在一个层级上的目标写成
\begin{equation}
 \begin{aligned}
  L_{\mathrm{dense}}
  &=\frac{1}{|M|}\sum_{i\in M}
    \lVert\widehat{\bm{h}}_i-\operatorname{sg}(\bm{h}^{*}_i)\rVert_1
   \\
  &\quad+\frac{1}{|C|}\sum_{i\in C}\lambda_i
    \lVert\widehat{\bm{h}}_i-\operatorname{sg}(\bm{h}^{*}_i)\rVert_1\text{。}
 \end{aligned}
  \label{eq:vision-ivm-vjepa-dense}
\end{equation}
$\operatorname{sg}$表示目标端停止梯度，即本次误差不直接更新教师。抽样要求$M$、$C$均非空；两个分母分别按各组位置数归一化，避免遮蔽比例决定哪组占据更多权重。原方法令$\lambda_i=\lambda/\sqrt{d_{\min}(i,M)}$，其中$d_{\min}$以块网格计量可见位置到最近遮蔽位置的距离，$\lambda$控制总体强度。靠近缺失区域的可见位置得到较大权重，用来加强局部连接。

多层监督的做法是：读取编码器三个中间层和最终层的输出，把各层归一化后的特征沿通道拼接，经小网络压回预测器所需的维度，再为四个层级分别预测教师的对应层特征，每层按式\eqref{eq:vision-ivm-vjepa-dense}计算误差并相加。训练从图像和视频中抽样，两种输入分别使用二维和三维切块投影，共享主体编码器。学生和预测器由梯度更新，教师按$\bar{\bm\theta}\leftarrow\mu\bar{\bm\theta}+(1-\mu)\bm\theta$缓慢跟随学生，$0<\mu<1$。这里的教师和卷八中的DINO教师一样提供表示目标，并非真实深度或动作标签。

各组件都是标准ViT。切块层按模态分开：视频用核与步长均为$2\times16\times16$的三维卷积，图像用$16\times16$的二维卷积；V-JEPA 2把图像在时间上复制成16帧的静止视频，V-JEPA 2.1不再这样做，而是在编码器和预测器的输入中各加一个可学习的模态词元，标明当前样本来自图像还是视频。按这一切块方式，主训练阶段$16$帧$256\times256$的视频有$8\times16\times16=2048$个位置，$256\times256$的图像有$16\times16=256$个位置。表\ref{tab:vision-ivm-vjepa-configs}列出编码器与预测器的配置。

\begin{table*}[tbp]
  \centering
  \small
  \begin{tabular}{@{}llrrrrr@{}}
    \toprule
    组件 & 模型 & 层数 & 宽度 & 头数 & MLP宽度 & 参数量\\
    \midrule
    编码器 & ViT-L & 24 & 1024 & 16 & 4096 & 3亿\\
    编码器 & ViT-g & 40 & 1408 & 22 & 6144 & 10亿\\
    编码器 & ViT-G & 48 & 论文未给出 & 论文未给出 & 论文未给出 & 20亿\\
    预测器 & 预训练用 & 24 & 384 & 论文未给出 & 论文未给出 & 论文未给出\\
    预测器 & 蒸馏用 & 12 & 384 & 论文未给出 & 论文未给出 & 论文未给出\\
    蒸馏学生 & ViT-B，ViT-L & — & — & — & — & 8000万，3亿\\
    \bottomrule
  \end{tabular}
  \caption{V-JEPA 2.1的编码器与预测器。ViT-L与ViT-g的层数、宽度、头数和MLP宽度取自V-JEPA 2论文的架构表，V-JEPA 2.1沿用这两个名称和参数规模；ViT-G的48层由论文给出的深监督层索引$[12,24,36,48]$确定，其余配置论文未给出。所有编码器的块大小为16，视频的时间块长为2。}
  \label{tab:vision-ivm-vjepa-configs}
\end{table*}

预测器比编码器窄得多，宽度只有384。它的任务是把可见位置的特征和遮蔽向量组成的序列读成每个位置的预测，层数从V-JEPA 2的12层增加到24层\scite{vision-vjepa2-2025}。编码器为ViT-G时，深监督读取第12、24、36、48个块的输出。

以单层的三个教学位置算一次误差：遮蔽位置的$L_1$残差为2，两个可见位置的残差分别为1和3，$\lambda=0.2$，到最近遮蔽位置的距离分别为1和4。两个可见权重为0.2和0.1，故本层目标为$2+(0.2\times1+0.1\times3)/2=2.25$。若教师某个参数旧值为1、更新后的学生值为1.2、$\mu=0.9$，教师新值为$0.9\times1+0.1\times1.2=1.02$。误差通过学生计算图更新预测，教师只按移动平均更新。

每次更新先抽取样本与空间增强、采样遮蔽区域，运行两端编码，累积各层误差，再反传并更新参数。论文的训练分两个阶段，各规模的编码器使用相同的超参数。主阶段共13.5万次迭代：视频为每秒4帧采样的16帧$256\times256$片段，图像为$256\times256$，每个全局批包含128段视频和2304张图像，学习率先在1.2万次迭代内从$10^{-4}$升到$5.25\times10^{-4}$再保持不变，教师的移动平均系数$\mu=0.99925$，可见位置损失的系数$\lambda$对视频取0.5、对图像取0.7。冷却阶段再训练1.2万次迭代，学习率降到$10^{-6}$，同时把视频改为64帧$384\times384$、图像改为$512\times512$；视频的位置数随之从2048增加到$32\times24\times24=18432$。蒸馏时把移动平均教师换成冻结的ViT-G，只在最后一层计算损失，不用深监督，并换用12层的新预测器。使用某个已发布权重时，要核对它来自哪个阶段、哪种分辨率和帧数。推理若只需要特征，保留编码器并输入完整观测，不必运行预测器和遮蔽。下游任务需要自己的监督：论文的深度评价冻结编码器，用最终层网格特征训练线性预测层；动作预期则冻结编码器，另训汇总特征的任务模块。论文报告，在NYUv2深度估计上以线性预测层得到0.307的均方根误差，在Something-Something V2动作识别上为77.7\%。其消融还显示，单独加入可见位置损失可改善密集任务却降低某些分类成绩，多层监督再改变这一取舍；密集监督的收益要按输出和训练配置核验。

\subsection{视频表示的形成与接口选择}
\label{sec:vision-ivm-video-interface}

固定拿杯视频的同一组请求，三个模型的差别可以归结到三个设计选择上。表\ref{tab:vision-ivm-video-models}列出这些差别。

\begin{table}[htbp]
  \centering
  \small
  \begin{tabularx}{\linewidth}{@{}p{16mm}XXX@{}}
    \toprule
    比较项 & TimeSformer & VideoMAE & V-JEPA 2.1\\
    \midrule
    输入单元 & 8帧逐帧图像块 & 16帧$2\times16\times16$时空块 & 图像与视频分别切块，共享编码器\\
    注意力结构 & 时间、空间分解 & 可见块间联合注意力 & 可见块间注意力\\
    监督来源 & 动作标签，图像预训练初始化 & 无标签视频，重建被遮像素 & 无标签图像与视频，预测教师表示\\
    受约束的位置 & 只有分类位置 & 被遮块 & 被遮块、可见块及多个中间层\\
    代表规模 & ViT-B结构，1.214亿参数 & ViT-B编码器，8700万参数 & ViT-L至ViT-G编码器，3亿至20亿参数\\
    主要代价 & 依赖标签，时间只在同位置比较 & 需拟合难预测的像素细节 & 需维护教师与预测器，结果依训练配置取舍\\
    \bottomrule
  \end{tabularx}
  \caption{同一拿杯视频请求下三个视频理解模型的比较。表中只列结构与训练上的差别；各自报告的数值来自不同设置，不能直接横向比较。}
  \label{tab:vision-ivm-video-models}
\end{table}
\FloatBarrier

VideoMAE改变了TimeSformer的监督来源：用管状掩码下的像素重建代替动作标签，换来无标签视频的预训练和只编码一成块的低预训练成本，代价是编码器要保留部分难以预测的像素细节，下游仍需标签微调。V-JEPA 2.1保留了VideoMAE的遮挡预测，改变的是预测目标和受约束的位置：预测教师表示而非像素，并让可见位置与中间层也受监督，换来冻结时即可用于深度等密集任务的特征，代价是训练中需要同时维护教师、预测器和多层目标，整体分类与密集预测之间的取舍也要按配置检验。

接口选择取决于要读的输出。只判断整段是否为“拿杯”时，三者都可以在全局向量上接分类头，TimeSformer在有充足标签时最直接；需要每帧中手的位置或远近时，V-JEPA 2.1的密集特征在冻结条件下更合适；目标领域有大量无标签视频而标签很少时，VideoMAE式的预训练可以先在领域数据上进行。三者都沿用了ViT的块序列和注意力，差别只在切块、注意力范围和训练目标上，这也使它们的权重可以互相作为初始化或比较的基准。\ref{chap:vision-video-recognition}章在具体任务中使用这些表示。

与\ref{chap:vision-video-foundations}章的经典方法相比，三个模型都没有显式的运动估计步骤。块匹配、Horn--Schunck和Lucas--Kanade输出每个位置的位移，后续处理可以直接读出“杯子向上移动了多少像素”；这三个模型输出的是时空位置上的特征，运动信息分散在注意力权重和特征之中，需要位移本身的任务仍要另接光流或跟踪模块，这些在\ref{chap:vision-video-recognition}章和\ref{chap:vision-video-restoration}章讨论。反过来，经典运动估计依赖亮度恒定和小位移等假设，遮挡和大幅运动时会失效；学习得到的视频表示不依赖这些显式假设，但它们对哪些运动敏感，只能由训练数据和评测结果说明。

\section{图像与视频生成模型}
\label{sec:vision-ivm-generation}

生成模型要输出新的图像或视频。直接在像素上生成代价很大：一张$256\times256\times3$的图像有约20万个数值，17帧$128\times128$的视频有约84万个，而扩散模型或自回归模型都要对输出反复计算许多步。常见的做法因此分两步：先训练一个编码器—解码器，把像素压缩成紧凑的表示并能还原；再在这个表示上训练生成模型。压缩表示必须同时满足两点：足够紧凑，使生成模型的计算可承受；保留足够的信息，使解码后的图像清晰可信。这一要求与\ref{sec:vision-transforms-multiresolution}节的变换编码相近，都是先换一种表示、再决定保留什么；不同之处在于，生成表示还要便于生成模型建模，而且可以用合成的方式而不是精确的逆变换还原像素，本节末尾再比较二者。

本节的三个模型对“生成在什么表示上进行、由什么主干完成”给出不同回答。VQGAN把图像压缩成离散码的网格，再用自回归Transformer逐个预测码；MAGVIT-v2仍用离散码，但改变了量化方式，使词表可以扩大到26万个码，并让图像和视频共用同一个编码器，再用掩码预测并行生成；DiT不使用离散码，而是在变分自编码器的连续潜表示上，用Transformer代替卷积网络做扩散去噪。前两者的差别在于离散表示怎样构造，后者的差别在于改用连续表示和扩散主干。卷六已经讲过变分自编码器和扩散模型的结构与训练，卷五讲过生成对抗网络及局部块判别器，本节在用到时只说明所需的结论。

\subsection{VQGAN}
\label{sec:vision-ivm-vqgan}

Transformer适合建模长距离的组合关系，例如街景中汽车、道路和天空的相对位置；但若把每个像素当作一个单元，$256\times256$的图像有65536个位置，注意力无法承受。VQGAN要解决的是：怎样把图像压成一个短得多的离散序列，使Transformer能在其上建模全局组合，同时解码后仍保留清晰的局部纹理\scite{vision-vqgan2021}。

它的主要想法是两阶段：第一阶段用卷积编码器—解码器学习一个由离散码组成的“图像部件词表”，卷积提供局部纹理的建模能力；第二阶段用Transformer学习这些码怎样组合成整幅图像。与之前的离散自编码器相比，VQGAN在第一阶段加入感知损失和局部块判别器，使得在较高压缩率下解码图像仍然清晰，从而把序列缩短到Transformer能处理的长度。

离散化依靠\term{向量量化}（vector quantization）：准备一个包含$K$个向量的\term{码本}（codebook）$\mathcal Z=\{\bm c_k\}_{k=1}^{K}$，把编码器输出的每个连续向量替换为码本中离它最近的那个。设编码器$E$把图像变成网格$\hat{\bm Z}=E(\bm X)\in\mathbb R^{h'\times w'\times d_z}$，对每个位置$(i,j)$，
\begin{equation}
  \bm z^{\mathrm q}_{i,j}=\bm c_{k^\ast},\qquad
  k^\ast=\argmin_{k}\lVert\hat{\bm z}_{i,j}-\bm c_k\rVert_2,
  \label{eq:vision-ivm-vq}
\end{equation}
解码器$G$再由量化网格还原图像$\hat{\bm X}=G(\bm Z^{\mathrm q})$。这样，一张图像就由$h'\times w'$个码的编号表示。这一离散自编码器的做法来自VQ-VAE\scite{vision-vqvae2017}。例如码本为$(0,0)$、$(1,0)$、$(0,1)$，编码器某位置输出$(0.8,0.3)$，它到三个码的距离约为0.85、0.36和1.06，于是该位置取第二个码。

取最近码是不可微的操作，梯度无法从解码器传回编码器。\term{直通估计}（straight-through estimator）的处理办法是：前向计算使用量化后的$\bm z^{\mathrm q}$，反向传播时把损失对$\bm z^{\mathrm q}$的梯度原样复制给$\hat{\bm z}$，好像量化这一步是恒等映射。码本本身则由另外的损失项训练。第一阶段的基本目标为
\begin{equation}
  L_{\mathrm{VQ}}=L_{\mathrm{rec}}(\bm X,\hat{\bm X})
  +\lVert\operatorname{sg}(\hat{\bm Z})-\bm Z^{\mathrm q}\rVert_2^2
  +\lVert\operatorname{sg}(\bm Z^{\mathrm q})-\hat{\bm Z}\rVert_2^2.
  \label{eq:vision-ivm-vq-loss}
\end{equation}
第二项只更新码本，把被选中的码拉向编码器输出；第三项称为承诺损失，只更新编码器，使编码器输出不偏离所选的码。没有第三项时，编码器输出可能越来越远离码本，量化误差随之增大。

VQGAN对重建项作了两处修改。其一，把逐像素的均方误差换成\term{感知损失}（perceptual loss），即比较两图在一个预训练网络中的特征差异，它不要求逐像素相同，对纹理和结构的偏差更敏感。其二，加入局部块判别器$D$，判断解码图像的每个局部块是否像真实图像块，编码器和解码器则要使判别器难以区分；对抗训练的机制可参考\BookChapterRef{chap:neural-generative-models}{卷五的神经生成模型章}。完整目标为
\begin{equation}
  \min_{E,G,\mathcal Z}\max_{D}\;
  \mathbb E\bigl[L_{\mathrm{VQ}}+\lambda L_{\mathrm{GAN}}\bigr],\qquad
  \lambda=\frac{\lVert\nabla_{G_L}L_{\mathrm{rec}}\rVert}
  {\lVert\nabla_{G_L}L_{\mathrm{GAN}}\rVert+\delta},
  \label{eq:vision-ivm-vqgan}
\end{equation}
其中$L_{\mathrm{GAN}}=\log D(\bm X)+\log(1-D(\hat{\bm X}))$，$G_L$是解码器最后一层，$\delta=10^{-6}$。$\lambda$按两项损失在最后一层上的梯度大小自动平衡：对抗项梯度过大时$\lambda$变小，避免它压过重建。训练初期先令$\lambda=0$，待重建稳定后再加入对抗项。

编码器和解码器的结构由论文附录表7给出，表\ref{tab:vision-ivm-vqgan-layers}按$256\times256$输入、下采样16倍的设置把它展开。编码器沿用DDPM去噪网络的卷积结构，但去掉了编码器与解码器之间的跳跃连接；若保留这些连接，未经量化的编码器特征会直接到达解码器，图像就不再只由码决定。下采样的次数$m$决定压缩倍数$f=2^m$，$f=16$对应$m=4$。每个阶段先经过残差块，再经过下采样块把高宽减半。最低分辨率处插入一个\term{非局部块}（non-local block），即在全部$16\times16$个位置之间做一次自注意力，让每个码在量化前就能读取整幅图像的上下文。编码器的最后一层把通道数变成码向量的维数$n_z$，量化后，解码器按相反的顺序上采样回像素。

\begin{table*}[tbp]
  \centering
  \small
  \begin{tabular}{@{}lllrl@{}}
    \toprule
    部分 & 层组 & 层类型 & 重复次数 & 输出尺寸\\
    \midrule
    编码器 & 输入 & 卷积 & 1 & $256\times256\times C'$\\
    编码器 & 下采样阶段 & 残差块，下采样块 & $m=4$ & $16\times16\times C''$\\
    编码器 & 中间层 & 残差块，非局部块，残差块 & 1 & $16\times16\times C''$\\
    编码器 & 输出 & 组归一化，Swish，卷积 & 1 & $16\times16\times n_z$\\
    量化 & 最近码替换 & 式\eqref{eq:vision-ivm-vq} & — & $16\times16$个编号\\
    解码器 & 输入 & 卷积 & 1 & $16\times16\times C''$\\
    解码器 & 中间层 & 残差块，非局部块，残差块 & 1 & $16\times16\times C''$\\
    解码器 & 上采样阶段 & 残差块，上采样块 & $m=4$ & $256\times256\times C'$\\
    解码器 & 输出 & 组归一化，Swish，卷积 & 1 & $256\times256\times3$\\
    \bottomrule
  \end{tabular}
  \caption{VQGAN编码器与解码器的分阶段结构，取自原论文附录表7，按$f=16$展开。论文没有给出各卷积的核大小以及通道数$C'$、$C''$；ImageNet实验中码向量维数$n_z=256$。判别器沿用pix2pix的局部块判别器，论文未另给层数。}
  \label{tab:vision-ivm-vqgan-layers}
\end{table*}

以论文的主要设置为例，码本大小$K=1024$，每个码向量256维，编码器把$256\times256$的图像下采样16倍，得到$16\times16$的码网格，即256个编号。原图需要$256\times256\times3\times8\approx157$万比特，256个编号每个只需10比特，共2560比特。第二阶段把这256个编号按光栅顺序排成序列$s$，用自回归Transformer建模
\begin{equation}
  p(s\mid c)=\prod_{i=1}^{256}p(s_i\mid s_{<i},c),
  \qquad L_{\mathrm{AR}}=\mathbb E\bigl[-\log p(s\mid c)\bigr],
  \label{eq:vision-ivm-vqgan-ar}
\end{equation}
其中$c$是条件：类别条件时为一个类别编号；分割图等空间条件时，先用另一个VQGAN把条件图编码为编号序列，放在$s$前面，只在$s$的位置计算损失。训练第二阶段时第一阶段的编码器、解码器和码本固定。生成时逐个采样256个编号，再经解码器得到图像；更高分辨率的图像则让Transformer在码网格上以滑动窗口逐块生成。这种逐个采样意味着256个编号需要256次顺序计算，这是自回归生成的主要代价。

第二阶段的Transformer采用GPT-2的解码器结构：每一层是带因果掩码的多头自注意力加前馈层，第$i$个位置只能读取它之前的编号和条件，最后一层在每个位置输出$K$类的softmax。论文所有实验都使用16个注意力头，表\ref{tab:vision-ivm-vqgan-transformer}列出类别条件ImageNet实验的两种规模。

\begin{table*}[tbp]
  \centering
  \small
  \begin{tabular}{@{}lrrrrrr@{}}
    \toprule
    配置 & 码本大小 & 序列长度 & 层数 & 宽度 & 头数 & 参数量\\
    \midrule
    类别条件ImageNet & 1024 & 257 & 24 & 1024 & 16 & 3.07亿\\
    类别条件ImageNet（大） & 16384 & 257 & 48 & 1536 & 16 & 14亿\\
    \bottomrule
  \end{tabular}
  \caption{VQGAN第二阶段Transformer的配置，取自原论文附录表8。序列长度257为1个类别条件加$16\times16=256$个图像编号，每个编号对应原图$16\times16$像素的区域；论文与先前自回归模型比较时使用较大的一组。论文未给出MLP宽度。}
  \label{tab:vision-ivm-vqgan-transformer}
\end{table*}

论文选择$16\times16$的码网格，是因为它是GPT-2 medium结构（3.07亿参数）在12GB显存的GPU上可以训练的最大序列长度。这说明压缩倍数$f$不只由重建质量决定，也由第二阶段能承受的序列长度决定：$f=8$时同一张图有1024个编号，注意力的位置对增加到16倍。

\subsection{MAGVIT-v2}
\label{sec:vision-ivm-magvit2}

VQGAN式的离散表示有两处限制。一是词表难以扩大：MAGVIT-v2的论文观察到，向量量化的码本变大时重建质量持续改善，但在这些码上训练的语言模型的生成质量先升后降，过大的词表反而有害。二是图像和视频各用各的编码器，生成视频的模型不能直接利用图像数据。MAGVIT-v2把离散编码器看作视觉生成的关键部件，同时处理这两个问题\scite{vision-magvit2-2024}。

它的主要想法有两点。第一，用\term{无查表量化}（lookup-free quantization, LFQ）代替码本：不再保存码向量，而是把编码器输出的每一维直接按符号量化为$\pm1$，码的编号就是这些符号组成的二进制数，词表大小由维数决定。第二，用因果三维卷积构成编码器，让第一帧单独编码，后续帧只依赖之前的帧，这样单张图像就是“只有第一帧的视频”，图像和视频可以共用一个编码器和同一个词表。

LFQ的定义如下。设编码器在某个位置的输出为$\bm z\in\mathbb R^{m}$，$m=\log_2K$，量化为
\begin{equation}
  q(z_j)=\operatorname{sign}(z_j)\in\{-1,1\},\qquad
  \operatorname{Index}(\bm z)=\sum_{j=1}^{m}2^{j-1}\,\mathbb I\{z_j>0\}.
  \label{eq:vision-ivm-lfq}
\end{equation}
例如$m=3$，$\bm z=(0.7,-0.2,1.3)$量化为$(1,-1,1)$，编号为$1+0+4=5$，词表大小为$2^3=8$。论文取$m=18$，词表大小为$2^{18}=262144$。由于不需要在码本中搜索最近邻，词表扩大不增加查表的计算。为了让各个码都被使用，训练中加入熵惩罚
\begin{equation}
  L_{\mathrm{ent}}=\mathbb E\bigl[H(q(\bm z))\bigr]-H\bigl[\mathbb E(q(\bm z))\bigr],
  \label{eq:vision-ivm-lfq-entropy}
\end{equation}
其中$H$是熵。第一项要求每个位置的分配是确定的，第二项要求整批数据的分配平均下来接近均匀，使各个码都被用到。这两个熵需要可微的分配概率，由编码输出与各码的接近程度给出；在各维独立的假设下，第一项可按维分解计算。论文报告，在LFQ下重建和生成质量都随词表增大而持续改善。

编码器的输入可以是图像，也可以是视频。\term{因果卷积}（causal convolution）在时间维只向过去填充$k_t-1$帧、不向未来填充，$k_t$是时间方向的核长，于是每一帧的特征只依赖它自己和之前的帧。时间方向以步长2做两次下采样，17帧依次变为9帧和5帧：设当前有$1+2t$帧，在前面补2帧后用核长3、步长2的卷积，输出$t+1$个时间位置，其中第一个位置只看到第一帧。两次合起来，$1+4t$帧被压成$1+t$个时间位置，第一帧单独占一个，图\ref{fig:vision-ivm-causal-temporal}画出这一过程。17帧$128\times128$的视频经空间8倍、时间4倍压缩，得到$5\times16\times16$个位置，每个位置一个18比特的编号，共1280个编号。单张图像就是$t=0$的情形，只占第一个时间位置，因此可以与视频共用编码器和词表。这种组织方式与\ref{sec:vision-videofound-formats}节的视频编码格式相似：编码器先独立编码一帧，后续帧只参考已经编码的帧。差别在于，传统编码器显式估计运动并编码残差，MAGVIT-v2的编码器没有运动估计步骤，时间上的冗余由因果卷积在训练中学会利用。

\begin{figure}[htbp]
  \centering
  % MAGVIT-v2 编码器的时间因果下采样（核长3、步长2、向过去补2帧），只画时间维；宽约100mm。
\begin{tikzpicture}[font=\figurefont\small,x=4.3mm,y=1cm,
  frame/.style={draw=FigureBlue,fill=FigureBlueFill,line width=0.6pt,
    minimum width=3.5mm,minimum height=3.5mm,inner sep=0pt},
  mid/.style={frame,draw=FigureYellow,fill=FigureYellowFill},
  top/.style={frame,draw=FigureRed,fill=FigureRedFill},
  pad/.style={frame,draw=BookMuted,fill=black!8},
  link/.style={draw=BookMuted,line width=0.5pt},
  note/.style={text=BookMuted,font=\figurefont\footnotesize,anchor=east}]
  % 第一行：17 帧输入，前补 2 帧
  \node[note] at (-3.7,0) {17 帧输入};
  \node[pad] at (-1,0) {};
  \node[pad] at (0,0) {};
  \foreach \i in {1,...,17}{\node[frame] (a\i) at (\i,0) {};}
  \node[font=\figurefont\footnotesize] at (1,0.45) {1};
  \node[font=\figurefont\footnotesize] at (17,0.45) {17};
  % 第二行：9 个位置，与各自窗口中最新的一帧对齐，前补 2 个位置
  \node[note] at (-3.7,-1.3) {9 个位置};
  \node[pad] at (-3,-1.3) {};
  \node[pad] at (-1,-1.3) {};
  \foreach \j in {1,...,9}{
    \pgfmathtruncatemacro{\x}{2*\j-1}
    \node[mid] (b\j) at (\x,-1.3) {};}
  % 第三行：5 个位置
  \node[note] at (-3.7,-2.6) {5 个位置};
  \foreach \k in {1,...,5}{
    \pgfmathtruncatemacro{\x}{4*\k-3}
    \node[top] (c\k) at (\x,-2.6) {};}
  % 示例窗口：第二行第 3 个位置读第 3、4、5 帧；第三行第 2 个位置读第二行第 1、2、3 个位置
  \foreach \i in {3,4,5}{\draw[link] (a\i.south) -- (b3.north);}
  \foreach \j in {1,2,3}{\draw[link] (b\j.south) -- (c2.north);}
  % 第一个位置只读补位与第 1 帧
  \draw[link,dashed] (-1,-0.18) -- (b1.north);
  \draw[link,dashed] (0,-0.18) -- (b1.north);
  \draw[link] (a1.south) -- (b1.north);
  \draw[link,dashed] (-3,-1.48) -- (c1.north);
  \draw[link,dashed] (-1,-1.48) -- (c1.north);
  \draw[link] (b1.south) -- (c1.north);
\end{tikzpicture}

  \caption{MAGVIT-v2编码器中时间方向的因果下采样。每一行是一个时间位置序列，灰色方格是向过去补的帧；每个输出位置由核长为3、步长为2的窗口计算，窗口只覆盖当前和更早的位置。17帧经两次下采样变为9个、5个位置，第一个位置始终只依赖第一帧。图中只画时间维，空间维同时按表\ref{tab:vision-ivm-magvit2-layers}下采样。}
  \label{fig:vision-ivm-causal-temporal}
\end{figure}

编码器、解码器和判别器的逐阶段结构由论文附录图7和超参数表给出，表\ref{tab:vision-ivm-magvit2-layers}按17帧$128\times128$的视频输入整理编码器部分。基础通道数为128，四个分辨率层级的通道倍数为1、2、2、4，即128、256、256、512；每个层级有4个残差块。残差块的内部顺序是：组归一化，Swish激活，$3\times3\times3$的时间因果卷积，再重复一次组归一化、Swish和卷积，最后与捷径相加；输入输出通道不同时，捷径为一个$1\times1\times1$卷积。

\begin{table*}[tbp]
  \centering
  \small
  \begin{tabularx}{\linewidth}{@{}llXrrl@{}}
    \toprule
    阶段或层组 & 层类型 & 卷积核与步长 & 输出通道 & 重复次数 & 输出尺寸\\
    \midrule
    输入 & — & — & 3 & — & $17\times128\times128$\\
    输入卷积 & 时间因果卷积 & $3\times3\times3$，步长1 & 128 & 1 & $17\times128\times128$\\
    层级1 & 残差块 & $3\times3\times3$ & 128 & 4 & $17\times128\times128$\\
    下采样1 & 时间因果卷积 & $3\times3\times3$，步长$1\times2\times2$ & 128 & 1 & $17\times64\times64$\\
    层级2 & 残差块，首块$128\to256$ & $3\times3\times3$ & 256 & 4 & $17\times64\times64$\\
    下采样2 & 时间因果卷积 & $3\times3\times3$，步长$2\times2\times2$ & 256 & 1 & $9\times32\times32$\\
    层级3 & 残差块 & $3\times3\times3$ & 256 & 4 & $9\times32\times32$\\
    下采样3 & 时间因果卷积 & $3\times3\times3$，步长$2\times2\times2$ & 256 & 1 & $5\times16\times16$\\
    层级4 & 残差块，首块$256\to512$ & $3\times3\times3$ & 512 & 4 & $5\times16\times16$\\
    中间 & 残差块 & $3\times3\times3$ & 512 & 4 & $5\times16\times16$\\
    输出 & 组归一化，Swish，卷积 & $1\times1\times1$ & 18 & 1 & $5\times16\times16$\\
    量化 & LFQ & 每维取符号 & 18位 & — & $5\times16\times16$个编号\\
    \bottomrule
  \end{tabularx}
  \caption{MAGVIT-v2视频分词器编码器的结构，按原论文附录图7整理，输入为17帧$128\times128$的视频。解码器与之对称：先用卷积把18维变回512通道，逐级用“卷积后深度到空间重排”的方式上采样，并在每个分辨率的残差块之前加一个以量化结果为条件的自适应组归一化。视频分词器由在ImageNet上训练的二维模型膨胀为三维后初始化。}
  \label{tab:vision-ivm-magvit2-layers}
\end{table*}

相对于前一代MAGVIT，论文对结构做了几处改动并逐项消融：编码器的下采样从平均池化改为步长卷积，解码器的上采样从最近邻放大加卷积改为\term{深度到空间重排}（depth-to-space），即先用卷积把通道扩大$2\times2\times2=8$倍，再把这些通道重新排成8倍的时空位置；时间下采样从编码器的前几个块推迟到最后几个块；判别器的下采样改用三维模糊池化，以增强平移不变性；解码器每个分辨率前加入自适应组归一化，把量化结果作为控制信号，做法来自StyleGAN。判别器由一个$3\times3\times3$卷积、五个带模糊池化的下采样残差块、一个卷积和两层全连接组成，通道倍数为2、4、4、4、4。

第一阶段的训练目标包括重建误差、对抗损失、感知损失、承诺损失、熵惩罚和稳定判别器的LeCAM正则项及$R_1$梯度惩罚，不再有码本损失，因为已经没有需要学习的码向量。论文给出的权重为：重建5.0，生成器对抗损失0.1，感知损失0.1，承诺损失0.25，熵惩罚0.1，LeCAM 0.001，$R_1$惩罚10。视频分词器在Kinetics-600上训练190个轮次，批大小256。图像生成另训练两个只处理图像的分词器，空间分别下采样16倍和32倍，用于$256\times256$和$512\times512$的生成，两者都把一张图像表示为$16\times16$个编号，在ImageNet上训练270个轮次。

生成阶段使用MAGVIT的掩码语言模型：训练时随机遮住一部分编号，Transformer根据其余编号预测被遮的编号，损失为被遮位置的交叉熵。这个Transformer沿用MAGVIT的结构，参数约3亿（论文表中为307M）。$2^{18}$个类别的输出层对它过大，论文因此把18位编号拆成两个9位的子编号：输入端把两个子编号的嵌入相加作为该位置的嵌入，输出端用两个各含$2^9=512$类的预测头分别预测，并让嵌入矩阵与输出层共享权重。训练时有10\%的批次丢弃类别标签，以便生成时使用无分类器引导。生成时从全部遮住开始，每一步预测所有被遮位置，保留置信度高的编号，把置信度低的重新遮住，逐步减少被遮的数量。与式\eqref{eq:vision-ivm-vqgan-ar}的逐个生成相比，每一步并行预测多个编号，步数可以远少于序列长度。论文报告，在类别条件的ImageNet $512\times512$生成上，使用引导、64步解码时FID为1.91，低于其比较的最好扩散模型的2.65；FID越低表示生成图像的特征分布越接近真实图像。论文还在视频压缩上与传统编码器比较：在每像素0.0384比特时，MAGVIT-v2的LPIPS为0.104，优于VVC的0.153，但PSNR为26.18，低于VVC的32.65。感知相似度好而逐像素误差大，说明这种编码器保留的是看起来相似的内容，而不是逐像素的原值。

\subsection{DiT}
\label{sec:vision-ivm-dit}

扩散模型从噪声出发，经许多步去噪得到样本；卷六给出了它的前向加噪、反向去噪和噪声回归目标，可参考\BookChapterRef{chap:pgm-generative}{卷六的概率生成模型章}。潜空间扩散模型把扩散放到自编码器的连续潜表示上，用带有下采样和上采样路径的卷积U-Net作为去噪网络，\ref{sec:vision-models-ldm}节将介绍这一组合。DiT要回答的问题是：去噪网络能否换成标准Transformer，以及它能否像图像理解中的ViT那样随计算量的增加而稳定地变好\scite{peebles2023dit}。

它的主要想法是把ViT的块序列结构用于潜表示上的去噪，并通过自适应层归一化把噪声步和类别条件注入每一层。这样，扩散模型的主干与ViT等模型共用同一套结构，模型规模可以按层数、宽度和块大小系统地调整。

DiT的输入是潜表示而非像素。$256\times256\times3$的图像经预训练变分自编码器的编码器压缩8倍，得到$32\times32\times4$的潜数组$\bm z$；按扩散过程加噪得到第$k$步的$\bm z^{(k)}$。DiT把$\bm z^{(k)}$切成$p\times p$的块，块数为$(32/p)^2$：$p=2$时为256个，$p=4$时为64个，$p=8$时为16个；块大小减半，块数变为4倍，计算量也大幅增加。每块经线性映射得到$d$维向量，加上ViT中使用的正弦余弦位置编码。噪声步$k$和类别$y$按ADM的方式各自嵌入为向量$\bm e_k$和$\bm e_y$。潜数组上的一个$2\times2$块对应原图$16\times16$像素，与ViT-B/16的一个图像块覆盖相同的范围；差别在于ViT的块里是原始像素，DiT的块里是自编码器已经压缩过的8倍下采样特征。论文指出，改变$p$几乎不改变参数量，却至少按块数成比例地改变计算量。

条件注入使用\term{自适应层归一化}（adaptive layer normalization, adaLN）：层归一化后的缩放和平移量不是固定参数，而是由条件计算得到。DiT采用的adaLN-Zero还为每个残差分支增加一个门控系数。设一层的输入为$\bm H$，由$\bm e_k+\bm e_y$经一个小网络回归出$(\bm\gamma_1,\bm\beta_1,\bm\alpha_1,\bm\gamma_2,\bm\beta_2,\bm\alpha_2)$，则
\begin{equation}
  \begin{aligned}
    \bm H'&=\bm H+\bm\alpha_1\odot\operatorname{Attn}\bigl(\bm\gamma_1\odot\operatorname{LN}(\bm H)+\bm\beta_1\bigr),\\
    \bm H''&=\bm H'+\bm\alpha_2\odot\operatorname{MLP}\bigl(\bm\gamma_2\odot\operatorname{LN}(\bm H')+\bm\beta_2\bigr),
  \end{aligned}
  \label{eq:vision-ivm-adaln}
\end{equation}
其中各系数按特征维广播到所有块，$\odot$为逐元素乘积。回归$\bm\alpha$的网络初始化为零，于是训练开始时$\bm\alpha_1=\bm\alpha_2=\bm 0$，每一层都是恒等映射，深层网络从一个稳定的起点开始训练。论文比较了把条件作为额外的序列位置、通过交叉注意力读取条件、普通adaLN和adaLN-Zero四种方式，adaLN-Zero的生成质量最好，增加的计算与普通adaLN相当，而交叉注意力方式要多出约15\%的计算。

把式\eqref{eq:vision-ivm-adaln}展开成一个DiT块内部的计算顺序：先由条件向量$\bm e_k+\bm e_y$经MLP回归出6组$d$维系数；序列做层归一化，按$\bm\gamma_1,\bm\beta_1$缩放平移，进入多头自注意力，输出乘以门控$\bm\alpha_1$后加回输入；再做层归一化，按$\bm\gamma_2,\bm\beta_2$缩放平移，进入逐位置的MLP，输出乘以$\bm\alpha_2$后加回。与式\eqref{eq:vision-ivm-vit-block}的ViT块相比，层归一化的缩放平移由条件决定，并多了两个门控系数；注意力和MLP本身不变。条件对所有块施加同一组系数，它能改变“整幅图朝哪个类别、哪个噪声水平去噪”，却不能对不同位置施加不同的条件，这是adaLN相对于交叉注意力的局限，按文字生成的模型因此常改用交叉注意力（\ref{chap:vision-multimodal-models}章）。

论文沿用ViT的配置给出S、B、L三个规模，另加一个XL规模，再与$p=2,4,8$三种块大小组合成12个模型。表\ref{tab:vision-ivm-dit-configs}列出它们在$256\times256$图像（$32\times32\times4$潜数组）上的形状和计算量。

\begin{table*}[tbp]
  \centering
  \small
  \begin{tabular}{@{}lrrrrrrrr@{}}
    \toprule
    模型 & 层数 & 宽度 & 头数 & MLP倍率 & 参数量 & $p=8$ & $p=4$ & $p=2$\\
    \midrule
    DiT-S & 12 & 384 & 6 & 4 & 3300万 & 0.36 & 1.41 & 6.06\\
    DiT-B & 12 & 768 & 12 & 4 & 1.3亿 & 1.42 & 5.56 & 23.01\\
    DiT-L & 24 & 1024 & 16 & 4 & 4.58亿 & 5.01 & 19.70 & 80.71\\
    DiT-XL & 28 & 1152 & 16 & 4 & 6.75亿 & 7.39 & 29.05 & 118.64\\
    \bottomrule
  \end{tabular}
  \caption{DiT的配置与每次前向的计算量（单位GFLOPs），输入为$256\times256$图像的$32\times32\times4$潜数组，块大小$p=8,4,2$分别对应16、64、256个词元。层数、宽度和头数取自原论文表1，参数量与三种块大小下的计算量取自论文附录的完整结果表；参数量随块大小只有不到100万的差别，表中取$p=2$的值。MLP倍率按ViT配置为4，论文未单独列出。$512\times512$图像的DiT-XL/2有1024个词元，计算量为524.6 GFLOPs。}
  \label{tab:vision-ivm-dit-configs}
\end{table*}

表中同一列从上到下、同一行从左到右，计算量都在增加。以DiT-XL为例，块大小从8减到2，词元数从16增到256，计算量从7.39增到118.64 GFLOPs，约16倍；参数量几乎不变。

最后一层把每个块的$d$维输出映射为$p\times p\times2C$个数，$C=4$是潜表示的通道数，再按位置排回$32\times32\times8$：前4个通道是预测的噪声，后4个通道是逐元素的方差。噪声预测用卷六中的噪声回归目标$L_{\mathrm{simple}}=\mathbb E\lVert\bm\epsilon-\bm\epsilon_{\bm\theta}(\bm z^{(k)},k,y)\rVert_2^2$训练，方差则用完整的变分下界训练。最后一层先做自适应层归一化，再用线性层输出；这个线性层初始化为零，其余参数按ViT的标准方法初始化。扩散过程沿用ADM的设置：$k_{\max}=1000$步，方差按线性日程从$10^{-4}$增到$2\times10^{-2}$。自编码器取自Stable Diffusion的预训练变分自编码器，训练DiT时保持冻结。数据为ImageNet，训练使用AdamW、常数学习率$10^{-4}$、批大小256，不用权重衰减和学习率预热，只做水平翻转增强，并保存衰减率0.9999的参数指数移动平均用于生成。生成时从高斯噪声出发，按卷六的反向过程逐步去噪，论文报告结果时使用250步DDPM采样，最后用变分自编码器的解码器把潜表示还原为图像。

论文最重要的观察是计算量与生成质量的关系。在不同层数、宽度和块大小的组合中，模型每次前向的计算量越大，FID越低，而计算量相近的配置即使参数量不同，FID也相近。最大的DiT-XL/2（块大小2）在$256\times256$的类别条件ImageNet上训练700万步后，使用无分类器引导时FID为2.27，不使用引导时为9.62。这一关系说明在DiT的设计范围内，增加计算是改善生成质量的可靠途径；它不说明任意结构都遵循同样的规律。

\subsection{生成表示的形成与接口选择}
\label{sec:vision-ivm-generation-interface}

三个模型都把生成拆成“压缩表示”和“在表示上生成”两步，差别在于这两步分别怎样做。表\ref{tab:vision-ivm-generation-models}在生成$256\times256$街景的同一请求下比较它们。

\begin{table}[htbp]
  \centering
  \small
  \begin{tabularx}{\linewidth}{@{}p{16mm}XXX@{}}
    \toprule
    比较项 & VQGAN & MAGVIT-v2 & DiT\\
    \midrule
    表示形式 & 离散码，$K=1024$ & 离散码，$K=2^{18}$，无码本 & 连续潜数组\\
    $256\times256$图像的表示 & $16\times16$个编号 & $16\times16$个编号 & $32\times32\times4$\\
    视频 & 不直接支持 & 因果三维编码，图像视频共用 & 原文只处理图像\\
    生成主干 & 自回归Transformer & 掩码预测Transformer & 扩散去噪Transformer\\
    生成步数 & 等于序列长度 & 远少于序列长度的并行步 & 由扩散采样步数决定\\
    生成主干规模 & 24层，3.07亿参数 & 约3亿参数 & DiT-XL/2，28层，6.75亿参数\\
    主要代价 & 词表难扩大，逐个采样慢 & 编码器训练目标多，像素保真有限 & 每步处理全部块，需多步去噪\\
    \bottomrule
  \end{tabularx}
  \caption{同一生成请求下三个生成模型的比较。各模型报告的FID来自不同设置，表中不作数值比较。}
  \label{tab:vision-ivm-generation-models}
\end{table}
\FloatBarrier

MAGVIT-v2改变的是VQGAN的量化方式和编码器结构：无查表量化使词表可以扩大而生成质量不降，因果三维卷积使图像和视频共用编码器，掩码预测又把逐个生成换成并行生成；代价是编码器的训练目标更复杂，所压缩的内容偏重感知相似而非逐像素保真。DiT则放弃离散码，改用连续潜表示和扩散去噪，并把去噪网络从卷积U-Net换成Transformer；换来的是与ViT相同的可扩展结构和随计算量稳定改善的生成质量，代价是每一步都要处理全部块，并需要多步去噪。

这三种生成表示也可以与\ref{sec:vision-transforms-multiresolution}节的变换编码对照。JPEG式的压缩把图像切成$8\times8$的块做离散余弦变换，再对系数逐个标量量化并做熵编码（\ref{sec:vision-transforms-dct}节）；小波编码则逐级把图像分解为低频近似和各方向的细节子带（\ref{sec:vision-transforms-wavelet}节）。这些变换都是固定的正交或双正交变换，在不量化时可以精确还原，压缩只发生在量化这一步，误差可以直接用均方误差控制和预测，失真表现为块效应、振铃和模糊。VQGAN和MAGVIT-v2的编码器是学习得到的非线性变换，没有逆变换可用，解码器要根据码重新合成像素，量化也从逐系数的标量量化变成对整个向量的向量量化或逐维取符号。按前文的数字算，VQGAN的$16\times16$个10比特编号相当于每像素$2560/65536\approx0.039$比特，与MAGVIT-v2视频压缩比较中所取的每像素0.0384比特同一量级；在这样的码率下，解码器靠对抗损失和感知损失补出看起来合理的纹理，失真的表现因此不是模糊，而是细节被替换。DiT所用的潜数组不做量化，$32\times32\times4=4096$个数相对原图的$196608$个数减少到约四十八分之一，它的目的不是降低码率，而是让扩散模型在较小的数组上计算。经典变换编码选择基的标准是能量集中和可逆，生成表示还多了一条标准：生成模型能否在其上建模。离散余弦变换满足前两条，但它的系数是否便于生成模型建模，不在它的设计目标之内。

这些表示和主干被\ref{chap:vision-multimodal-models}章的多模态生成模型组合使用。LDM（\ref{sec:vision-models-ldm}节）的自编码器与VQGAN同样使用感知损失和局部块判别器训练，并提供带向量量化和带KL正则的两种潜空间，去噪网络则是卷积U-Net。Wan（\ref{sec:vision-models-wan}节）的视频自编码器与MAGVIT-v2一样让第一帧单独编码、后续帧按组因果压缩，但输出连续潜表示；它在这一潜空间上用Transformer主干按流匹配目标生成视频，延续了DiT的思路。LTX-2（\ref{sec:vision-models-ltx2}节）为视频和音频分别准备因果变分自编码器，再让两个生成分支交换信息；Show-o2（\ref{sec:vision-models-showo2}节）则在三维因果变分自编码器的潜空间上同时提取理解和生成所需的特征。选择生成表示时，要依次确认：需要图像还是视频，生成主干要求离散编号还是连续数组，解码后保留的是逐像素内容还是感知相似的内容。

\section{模型的使用约束}
\label{sec:vision-ivm-constraints}

前三节的模型形成了不同的表示，使用时都要回答同一个问题：这个表示能否支持当前任务所需的结果。模型给出的特征或生成的画面本身不会说明它缺了什么，约束需要从模型的计算过程推出。可以从三个方面检查：表示在编码和采样中保留了哪些信息，训练数据决定了表示覆盖哪些内容，计算资源决定了能处理多大的输入和多少候选。三个方面相互影响，例如提高分辨率可以保留更多信息，却增加计算；扩大训练数据可以扩展覆盖范围，也要求更多训练资源。

\subsection{信息保留的约束}
\label{sec:vision-ivm-information}

编码中每一次降低分辨率都会丢掉一部分细节。ResNet最后一个阶段的格子对应$32\times32$像素，ViT和DINOv3的块对应$16\times16$像素；街景中宽20像素的路牌在前者中不到一格，在后者中只占一两格。块内的纹理并没有全部消失，它仍可以写进几百个通道之中，但它在块内的具体位置不再由网格坐标给出，路牌上的文字因此难以从这些特征中读出。这与\ref{sec:vision-formation-sampling}节的采样问题同源：输出网格的间隔就是新的采样间隔，比它更细的结构无法再按位置分开。稠密预测头的上采样不能补回编码时丢失的信息，需要细节的任务只能提高输入分辨率、裁取局部区域重新编码，或使用保留更细尺度的特征。

视频还在时间上丢信息。TimeSformer以1/32的采样率取8帧，在每秒25帧的视频中相邻输入帧相隔约1.3秒，持续不到一秒的动作可能整段落在两帧之间；VideoMAE在Kinetics-400上每4帧取一帧、共16帧，覆盖64帧原始画面，时间更密而跨度更短。V-JEPA 2.1的时空块把两帧合成一个单元，单元的时间范围也随之变成一个区间。这与\ref{sec:vision-videofound-frame-rate}节讨论的时间采样是同一个问题：模型的有效帧率是原始帧率除以取帧步长，事件持续时间短于取帧间隔时，就可能整段不出现在输入中。对时间敏感的任务要记录这些采样设置，并检查目标事件的长度是否落在采样能看见的范围内。

生成表示的压缩率决定了解码后能还原什么。VQGAN把$256\times256$的图像压成256个10比特的编号，MAGVIT-v2的视频编码在感知相似上优于传统编码器，却在PSNR上落后；这说明这些表示保留的是看起来合理的纹理和结构，不保证逐像素还原。用这些表示做复原或压缩时，招牌上的文字和车牌号码可能被解码成形状相似但内容不同的字符，这种结果在第\ref{chap:vision-overview}章的意义上没有完成保真复原。

\subsection{训练数据的约束}
\label{sec:vision-ivm-data}

表示能区分什么，取决于训练数据和目标要求它区分什么。ResNet和ViT的类别监督只覆盖预训练的类别表，ImageNet中没有区分的细粒度差别，例如两种外观相近的车型，不一定在特征中得到保留。ViT对预训练数据量更敏感，只在中等规模数据上训练时可能不如卷积主干，原论文的对比已经说明了这一点。

自监督模型摆脱了类别标签，却仍受图像来源影响。DINOv3的LVD-1689M从公开网络图像中筛选而来，这些网络照片的视角、光照和内容分布决定了冻结特征在哪些领域可靠；医学影像、遥感图像或工业零件与这些图像差别较大，需要在目标领域上单独检验，必要时在领域数据上继续训练。VideoMAE的论文也报告，预训练数据与目标数据之间的领域差别会影响迁移，数据质量比数量更重要。

生成模型只能生成训练分布能够支持的内容。在类别条件的ImageNet上训练的DiT只接受一千个类别作为条件，无法按文字描述生成；按文字生成需要图文配对数据和相应的条件接口，这由\ref{chap:vision-multimodal-models}章说明。训练数据中的偏差也会进入生成结果，例如某类场景总在白天出现，模型生成夜景时就缺少可靠的依据。

\subsection{计算资源的约束}
\label{sec:vision-ivm-compute}

注意力的计算按单元数的平方增长，这是基于Transformer的视觉模型最主要的资源约束。$224\times224$的照片取$p=16$有196个块，每层约3.8万个位置对；分辨率提高到$448\times448$，块数变为784，位置对增加到约61万，是原来的16倍。16帧视频的1568个时空块若做联合注意力，每层约246万个位置对。TimeSformer的分解注意力、VideoMAE预训练时只编码一成可见块，都是针对这一增长的设计；但VideoMAE微调和推理时仍要处理全部1568个块。按论文给出的数字，ResNet-50处理一张$224\times224$的照片约需$3.8\times10^9$次乘加，TimeSformer处理一个8帧视图约需0.59 TFLOPs，VideoMAE的ViT-B处理一个16帧视图约需180 GFLOPs；各论文对计算量的计法不尽相同，这些数只用于判断量级。

模型规模决定训练和部署的代价。DINOv3约67亿参数的ViT-7B需要大规模计算资源训练，实际使用多取蒸馏得到的较小模型；蒸馏后的模型推理更快，但在需要细致局部特征的任务上，与原模型的差距需要按任务核验。TimeSformer的论文还给出了训练成本的对比：在Kinetics-400上达到约75.8\%的准确率，它用了约416个GPU小时，其比较的SlowFast达到相近准确率用了约3840个GPU小时；这一对比依赖其ImageNet初始化和具体硬件，不是所有视频任务上的一般结论。

生成模型还要考虑生成步数。VQGAN的自回归生成需要与序列长度相同的顺序步数，MAGVIT-v2的掩码预测用64步左右并行生成，DiT则需要扩散采样的多步去噪，每步都处理全部块。DiT的结果显示，增加每步的计算可以稳定改善质量，但这意味着生成一张图像的总计算等于每步计算乘以步数：DiT-XL/2每步约118.6 GFLOPs，按论文评测时的250步DDPM采样，生成一张$256\times256$的图像约需$3\times10^4$ GFLOPs，即约30 TFLOPs，还未计入无分类器引导所需的额外前向。选择模型时，应在任务允许的等待时间和可用硬件下，比较达到所需质量的总计算，而不是只比较参数量或单步速度。

\section{本章小结}
\label{sec:vision-ivm-summary}

开篇的街景和拿杯视频提出了一个问题：同样的像素，为什么分类、分割、深度和生成需要不同的表示，这些表示又由什么计算形成。本章的回答分两层。在计算流程上，编码决定了网络以什么单元计算：卷积网络在像素网格上逐阶段降低分辨率，得到多尺度特征图；Transformer式编码器把图像切成图像块、把视频切成时空块，单元数由块大小决定，位置和时间必须显式编码。解码决定了读出什么：全局向量适合分类和检索，空间特征图保留粗粒度位置，稠密预测头结合多个尺度恢复到像素分辨率，但不能补回编码时已经丢失的细节。

在模型上，三组模型分别回答了三个问题。图像理解中，ResNet依靠卷积结构和类别标签形成多尺度特征，ViT以全局注意力换取更大的数据需求，DINOv3在ViT之上改用自监督目标并以Gram锚定保持块特征的空间结构。视频理解中，TimeSformer分解时空注意力以降低计算，VideoMAE以管状掩码和像素重建摆脱标签，V-JEPA 2.1把预测目标换成教师表示并约束可见位置和中间层，使特征适合密集预测。生成中，VQGAN以码本和自回归生成离散码，MAGVIT-v2以无查表量化扩大词表并让图像和视频共用编码器，DiT在连续潜表示上以Transformer完成扩散去噪。每组模型依次改变了一个设计选择，换来某种能力，也付出了数据、计算或保真度上的代价。

这些模型的具体结构也说明了它们与第一部分经典方法的关系。ResNet-50的茎和四个阶段把$224\times224$的照片逐级降到$7\times7$，每个阶段由滤波和降采样组成，相当于一座由数据决定滤波器的图像金字塔，但步长卷积不保证降采样前的低通；ViT、TimeSformer、VideoMAE和DiT都先把图像、视频或潜数组切成固定大小的块，对每块做一次线性变换，与分块变换编码的组织相同，只是变换由训练决定、不要求可逆；视频模型取帧受时间采样的限制，却都没有显式的运动估计；VQGAN和MAGVIT-v2的分词器与变换编码一样先变换再量化，但解码器靠合成而非逆变换还原像素，因而保留的是感知上的相似。理解这些对应关系，可以帮助判断一个模型会在哪里丢失信息。

这些模型仍有本章没有解决的部分。它们只处理视觉输入，用户的文字请求、视频中的声音怎样与视觉表示对应，留给\ref{chap:vision-multimodal-models}章；表示中保留的细节、训练数据覆盖的范围和可承受的计算，也只给出了检查的方向，具体任务中的取舍要在后续三个部分的任务章中，结合监督、输出恢复和评价协议逐一确定。

\paragraph{扩展阅读}
图像主干方面，EfficientNet用一个复合系数同时放大网络的深度、宽度和输入分辨率，回答的是ResNet式卷积网络在给定计算预算下应当沿哪一维扩大\scite{vision-1905-11946}。Swin Transformer从ViT一侧向卷积网络靠拢：注意力只在不重叠的局部窗口内计算，相邻层平移窗口以交换信息，并逐阶段合并图像块，得到与ResNet相同的多尺度特征图，计算量随图像块数线性增长，可直接替换检测和分割中的卷积主干\scite{vision-2103-14030}。ConvNeXt则从ResNet出发，参照这类分层Transformer逐项改动块设计和训练配方，说明本章ResNet与ViT之间的一部分差距来自设计细节和训练设置，而不只来自卷积与注意力之分；DINOv3蒸馏出的卷积学生采用的就是这一结构\scite{vision-2201-03545}。

视频方面，I3D把在图像上预训练的二维卷积核沿时间“膨胀”为三维卷积核，用三维卷积直接处理时空块，同时发布了本章反复用作基准的Kinetics数据集\scite{vision-i3d2017}。SlowFast用低帧率、通道多的慢路径读取空间语义，用高帧率、通道少的快路径捕捉运动，两条路径横向连接；它是TimeSformer训练成本比较中的卷积对照\scite{vision-1812-03982}。ViViT与TimeSformer同期，系统比较了联合时空注意力、先空间后时间的分解编码器以及在层内分解注意力的几种方案，并提出把相邻帧同一位置的块合成一个时空块的嵌入方式，可与TimeSformer的分解注意力和VideoMAE的时空块对照阅读\scite{vision-2103-15691}。

生成方面，MAGVIT-v2所用的掩码并行生成来自MaskGIT：它沿用VQGAN式的离散码，训练双向Transformer预测被掩去的码，生成时每轮保留置信度最高的一部分，用十步左右的迭代代替逐个码的自回归解码\scite{vision-2202-04200}。Stable Diffusion 3在DiT之后把训练目标换成沿噪声与数据之间直线路径的整流流，并以MMDiT让文字词元和图像词元各用一套权重、在同一注意力中联合计算，还考察了这类模型随参数规模扩大时的表现\scite{vision-2403-03206}。
